
\newpage

\section{Notations for Tensors and Derivatives} \label{apx:tensor}
We follow the notations of tensors used in \cite{kolda2009tensor}.
For a three-way tensor $\fX \in \BR^{d_1 \times d_2 \times d_3}$,
we use $ [\fX]_{i_1,i_2,i_3}$ to represent its $(i_1,i_2,i_3)$-th element.
The inner product of two three-way tensors $\fX,\fY$ is defined by $\langle \fX, \fY \rangle := \sum_{i_1,i_2,i_3} [\fX]_{i_1,i_2,i_3} \cdot [\fY]_{i_1,i_2,i_3}$. The operator norm of three-way tensor $\fX \in \BR^{d_1 \times d_2 \times d_3}$ is defined by $\Vert \fX \Vert:= \sup_{\Vert x_i \Vert =1} \langle \fX, x_1 \circ x_2 \circ x_3 \rangle $, where the elements of  $x_1 \circ x_2 \circ x_3 \in \BR^{d_1 \times d_2 \times d_3}$ is $ [x_1 \circ x_2 \circ x_3]_{i_1,i_2,i_3} := [x_1]_{i_1} \cdot  [x_2]_{i_2} \cdot [x_3]_{i_3} $. 
It can be verified that such a definition generalizes the Euclidean norm of a vector and the spectral norm of a matrix to tensors.
For a three-way tensor  $ \fX \in \BR^{d_1 \times d_2 \times d_3}$ and a vector $v \in \BR^{d_1}$, their mode-1 product, denoted by $\fX \bar \times_1 v$, gives a matrix in $\BR^{d_2 \times d_3}$ with elements $ [\fX \bar \times_1 v]_{i_2,i_3} := \sum_{i_1} [\fX]_{i_1,i_2,i_3} \cdot [v]_{i_1} $.
We define $\bar \times_2$ and $\bar \times_3$ in a similar way, and it can  be verified that 
$ \Vert \fX \bar \times_i v \Vert \le \Vert \fX \Vert \Vert v \Vert$.
For a three-way tensor $\fX \in \BR^{d_1 \times d_2 \times d_3}$ and a matrix $A \in \BR^{d_1' \times d_1}$, their mode-1 product, denoted by $ \fX \times_1 A$, gives a tensor in $\BR^{d_1' \times d_2 \times d_3}$ with elements  $ [\fX \times_1 A]_{i_1',i_2,i_3} := \sum_{i_1} [\fX]_{i_1,i_2,i_3} \cdot [A]_{i_1',i_1} $. We define $ \times_2$ and $ \times_3$ in a similar way, and it can also be verified that 
$ \Vert \fX \times_i A \Vert \le \Vert \fX \Vert \Vert A \Vert$.

For a function $h(x,y) : \BR^{d_x} \times \BR^{d_y} \rightarrow \BR$,
we denote $\nabla_x h(x,y) \in \BR^{d_x}$ to be the partial gradient with respect to $x$, with elements given by $ [\nabla_x h(x,y) ]_{i} := \partial f(x,y) / \partial [x]_{i}$. And we define $\nabla_y h(x,y) \in \BR^{d_y}$ in a similar way.  We denote $ \nabla_{xx}^2 h(x,y) \in \BR^{d_x} \times \BR^{d_x}$ to be the partial Hessian with respect to $x$, with elements given by $ [\nabla_{xx}^2 h(x,y)]_{i_1,i_2} := \partial^2 h(x,y) / (\partial [x]_{i_1} \partial [x]_{i_2})$. And we define $ \nabla_{xy}^2 h(x,y) \in \BR^{d_x} \times \BR^{d_y} $, $\nabla_{yx}^2 h(x,y) \in \BR^{d_y} \times \BR^{d_x}$ and  $ \nabla_{yy}^2 h(x,y) \in \BR^{d_y} \times \BR^{d_y}$ in a similar way.  We denote $\nabla_{xxx}^3 h(x,y) \in \BR^{d_x} \times \BR^{d_x} \times \BR^{d_x}$ to be the partial third-order derivative with respect to $x$, with elements given by 
$ [\nabla_{xxx}^3 h(x,y)]_{i_1,i_2,i_3} := \partial^3 h(x,y) / (\partial [x]_{i_1} \partial [x]_{i_2} \partial [x]_{i_3})$. And we define $\nabla_{xxy}^3 h(x,y) \in \BR^{d_x} \times \BR^{d_x} \times \BR^{d_y}$, $\nabla_{xyx}^3 h(x,y) \in \BR^{d_x} \times \BR^{d_y} \times \BR^{d_x}$,
$\nabla_{xyy}^3 h(x,y) \in \BR^{d_x} \times \BR^{d_y} \times \BR^{d_y}$,
$\nabla_{yyy}^3 h(x,y) \in \BR^{d_y} \times \BR^{d_y} \times \BR^{d_y}$,
$\nabla_{yyx}^3 h(x,y) \in \BR^{d_y} \times \BR^{d_y} \times \BR^{d_x}$,
$\nabla_{yxy}^3 h(x,y) \in \BR^{d_y} \times \BR^{d_x} \times \BR^{d_y}$
and$\nabla_{yxx}^3 h(x,y) \in \BR^{d_y} \times \BR^{d_x} \times \BR^{d_x}$
in a similar way.
We denote $\nabla y^*(x) \in \BR^{d_x} \times \BR^{d_y} $ to be the matrix with elements given by $[\nabla y^*(x) ]_{i_1,i_2} :=  \partial [y^*(x)]_{i_2}/\partial [x]_{i_1}$. We denote $\nabla^2 y^*(x) \in \BR^{d_x} \times \BR^{d_x} \times \BR^{d_y} $ to be the three-way tensor with elements given by $[\nabla^2 y^*(x)]_{i_1,i_2,i_3}: = \partial^2 [y^*(x)]_{i_3} / (\partial [x]_{i_1} \partial [x]_{i_2})$. And we define $\nabla y_{\lambda}^*(x)$ and $\nabla^2 y_{\lambda}^*(x)$ in a similar way.


\section{Lemmas for Finding First-Order Stationarity} \label{apx:lem-1st}

\begin{lem}[\citet{kwon2023fully}] \label{lem:Lag-sc}
Under Assumption \ref{asm:sc}, for $\lambda \ge 2L_f /\mu$,   $\fL_{\lambda}(x,\,\cdot\,)$ is  $(\lambda \mu/2)$-strongly convex.  
\end{lem}
% \begin{proof}
% It can be verified that
% \begin{align*}
%     \nabla_{yy}^2 \fL(x,y) = \nabla_{yy}^2 f(x,y) +  \lambda \nabla_{yy}^2 g(x,y) \succeq \frac{\lambda \mu}{2} I_{d_y}.
% \end{align*}
% % Since  $g^*(x)$ is $(1+ L_g / \mu)L_g$-smooth, then the smoothness constant of $\fL_{\lambda}(x,y)$ can be bounded by:
% % \begin{align*}
% %     L_f + \lambda \left(2+ \frac{L_g}{\mu} \right)L_g \le \frac{4 \lambda L_g^2}{\mu^2}.
% % \end{align*}
% \end{proof}


% Then it is shown in \cite{kwon2023fully} that the following results hold. We provide the proof here for completeness.

\begin{lem} \label{lem:yy-lambda}
    Under Assumption \ref{asm:sc}, for $\lambda \ge 2L_f/ \mu$, it holds that 
    \begin{align*}
        \Vert y_{\lambda}^*(x) - y^*(x) \Vert \le \frac{C_f}{\lambda \mu}
    \end{align*}
\end{lem}

\begin{proof}
    % It is a special case of Lemma 3.2 in \cite{kwon2023fully} with $x_1 = x_2, \lambda_1 = \lambda $ and $\lambda_2 = + \infty $.
By the first-order optimality condition, we know that
\begin{align*}
    \nabla_y f(x, y_{\lambda}^*(x)) + \lambda \nabla_y g(x,y_{\lambda}^*(x))  = 0.
\end{align*}
Then we have
\begin{align*}
    \Vert y_{\lambda}^*(x) - y^*(x) \Vert 
    \le \frac{1}{\mu} \Vert \nabla_y g(x,y_{\lambda}^*(x)) \Vert = \frac{1}{\lambda \mu} \Vert \nabla_y f(x,y_{\lambda}^*(x)) \Vert \le \frac{C_f}{\lambda \mu}
\end{align*}
\end{proof}

As a direct consequence, we can show that $\fL_{x}^*$ and $\varphi(x)$ are close.
\begin{lem} \label{lem:F-Lag-close}
 Under Assumption \ref{asm:sc}, for $\lambda \ge 2L_f / \mu$, it holds that $\vert \fL_{\lambda}^*(x) - \varphi(x) \vert \le D_0/ \lambda$, where
 \begin{align*}
     D_0 : = \left(C_f + \frac{C_f L_g}{2 \mu} \right) \frac{C_f }{\mu} = \fO(\ell \kappa^2).
 \end{align*}
\end{lem}

\begin{proof}
A simple calculus shows that
    \begin{align*}
    &\quad \vert \fL_{\lambda}^*(x) - \varphi(x) \vert \\
    &\le   \vert f(x,y_{\lambda}^*(x)) - f(x,y^*(x)) \vert + \lambda \vert g(x,y_{\lambda}^*(x)) - g(x,y^*(x)) \vert \\
    &\le C_f  \Vert y_{\lambda}^*(x) - y^*(x) \Vert + \frac{\lambda L_g}{2} \Vert y_{\lambda}^*(x) - y^*(x) \Vert^2 \\
    &\le \left(C_f + \frac{C_f L_g}{2 \mu} \right) \Vert y_{\lambda}^*(x) - y^*(x) \Vert \\
    &\le \left(C_f + \frac{C_f L_g}{2 \mu} \right) \frac{C_f }{\lambda \mu}.
    \end{align*}
\end{proof}
The following result is the key to designing fully first-order methods for bilevel optimization, first proved in Lemma 3.1 in \cite{kwon2023fully}. We provide the proof here for completeness.
\begin{lem} \label{lem:key-fully}
Under Assumption \ref{asm:sc}, for $\lambda \ge 2L_f/ \mu$,  it holds that $\Vert \nabla \fL_{\lambda}^*(x) - \nabla \varphi(x) \Vert \le D_1 / \lambda$, where
\begin{align} \label{eq:D-lambda}
    D_1 :=  \left( L_f 
   + \frac{\rho_g L_g}{\mu} + \frac{C_f L_g \rho_g}{2\mu^2} + \frac{C_f \rho_g}{2\mu} 
    \right) \frac{C_f}{\mu} = \fO(\ell \kappa^3).
\end{align}
\end{lem}

\begin{proof}
Taking total derivative on $\varphi(x) = f(x,y^*(x))$, we obtain the following result:
\begin{align} \label{eq:hyper-grad}
    \nabla \varphi(x) = \nabla_x f(x,y^*(x)) - \nabla_{xy}^2 g(x,y^*(x)) [ \nabla_{yy} g(x,y^*(x))]^{-1} \nabla_y f(x,y^*(x)). 
\end{align}
Also, we know that
\begin{align*}
    \nabla \fL_{\lambda}^*(x) = \nabla_x f(x,y_{\lambda}^*(x)) + \lambda (\nabla_x g(x,y_{\lambda}^*(x)) - \nabla_x g(x,y^*(x))).
\end{align*}
By simple calculus, we have
\begin{align*}
    &\quad \nabla \fL_{\lambda}^*(x) - \nabla \varphi(x) \\
    &= \nabla_x f(x,y_{\lambda}^*(x)) - \nabla_x f(x,y^*(x)) \\
    &\quad + \nabla_{xy}^2 g(x,y^*(x)) [ \nabla_{yy} g(x,y^*(x))]^{-1} (\nabla_y f(x,y^*(x)) - \nabla_y f(x,y_{\lambda}^*(x)) \\
    &\quad + \lambda \nabla_{xy}^2 g(x,y^*(x)) [ \nabla_{yy} g(x,y^*(x))]^{-1} \times \\
    &~~~~~~~~~~~~~~~( \nabla_{yy}^2 g(x,y^*(x)) (y_{\lambda}^*(x) - y^*(x)) + \nabla_y g(x,y^*(x)) - \nabla_y g(x,y_{\lambda}^*(x)) ) \\
    &\quad + \lambda ( \nabla_x g(x,y_{\lambda}^*(x)) - \nabla_x  g(x,y^*(x)) - \nabla_{xy}^2 g(x,y^*(x)) (y_{\lambda}^*(x) - y^*(x)) ).
\end{align*}
Taking norm on both sides,
\begin{align*}
    \Vert \nabla \fL_{\lambda}^*(x) - \nabla \varphi(x) \Vert &\le L_f \Vert y_{\lambda}^*(x) - y^*(x) \Vert + \frac{\rho_g L_g}{\mu} \Vert y_{\lambda}^*(x) - y^*(x) \Vert  \\
    &\quad + \frac{\lambda L_g \rho_g}{2\mu} \Vert y_{\lambda}^*(x) - y^*(x) \Vert^2  + \frac{\lambda \rho_g}{2} \Vert y_{\lambda}^*(x) - y^*(x) \Vert^2 \\
    &\le \left( L_f 
   + \frac{\rho_g L_g}{\mu} + \frac{C_f L_g \rho_g}{2\mu^2} + \frac{C_f \rho_g}{2\mu} 
    \right) \Vert y_{\lambda}^*(x) - y^*(x) \Vert \\
    &\le \left( L_f 
   + \frac{\rho_g L_g}{\mu} + \frac{C_f L_g \rho_g}{2\mu^2} + \frac{C_f \rho_g}{2\mu} 
    \right) \frac{C_f}{\lambda \mu}
\end{align*}


\end{proof}

\begin{lem} \label{lem:yxyx-lambda}
Under Assumption \ref{asm:sc}, for $\lambda \ge 2 L_f / \mu$, it holds that $ \Vert \nabla y^*(x) - \nabla y_{\lambda}^*(x) \Vert \le D_2/\lambda $, where
\begin{align*}
    D_2 := \left( \frac{1}{\mu} + \frac{2L_g}{\mu^2} \right) \left( L_f + \frac{C_f \rho_g}{\mu} \right) = \fO(\kappa^3).
\end{align*}

\end{lem}


\begin{proof}
 Taking derivative on both sides of $\nabla_y g (x,y^*(x)) = 0$ yields
 \begin{align} \label{eq:yxyy}
     \nabla_{xy}^2 g(x,y^*(x)) +  \nabla y^*(x)\nabla_{yy}^2 g(x,y^*(x)) = 0.
 \end{align}
 Rearranging, we can obtain
\begin{align} \label{eq:y-x}
  \nabla y^*(x) = - \nabla_{xy}^2 g(x,y^*(x)) [\nabla_{yy}^2 g(x,y^*(x)) ]^{-1} . 
  \end{align}
Similarly, we also have
\begin{align} \label{eq:y-lambda-x}
    \nabla y_{\lambda}^*(x) = -  \nabla_{xy}^2 \fL_{\lambda}(x,y_{\lambda}^*(x)) [\nabla_{yy}^2 \fL_{\lambda}(x,y_{\lambda}^*(x)) ]^{-1}.
\end{align}
Using the matrix identity $ A^{-1} - B^{-1} = A^{-1} (B -A) B^{-1}$, we have
\begin{align} \label{eq:diff-inv}
    \begin{split}
    &\quad  \left \Vert [\nabla_{yy}^2 g(x,y^*(x))]^{-1} - \left [  \frac{\nabla_{yy}^2 \fL_{\lambda}(x,y_{\lambda}^*(x))}{\lambda} \right]^{-1}  \right \Vert \\
    &\le  \left \Vert [\nabla_{yy}^2 g(x,y_{\lambda}^*(x))]^{-1} \right \Vert  
    \left \Vert  \frac{\nabla_{yy}^2 \fL_{\lambda}(x,y_{\lambda}^*(x))}{\lambda} -   \nabla_{yy}^2 g(x,y^*(x))  \right \Vert \left \Vert \left[\frac{\nabla_{yy}^2 \fL_{\lambda} (x,y_{\lambda}^*(x))}{\lambda}\right]^{-1}  \right \Vert \\
    &\le \frac{2}{\mu^2} \left \Vert 
    \frac{\nabla_{yy}^2 f(x,y_{\lambda}^*(x))}{\lambda}
     + \nabla_{yy}^2 g(x,y_{\lambda}^*(x)) - \nabla_{yy}^2 g(x,y^*(x)) \right \Vert \\
     &\le \frac{2}{\mu^2} \left( \frac{L_f}{\lambda} + \rho_g \Vert y_{\lambda}^*(x) - y^*(x) \Vert  \right)\\
     &\le \frac{2}{\lambda \mu^2} \left(L_f + 
     \frac{C_f \rho_g}{\mu}
     \right). 
    \end{split}
\end{align}
Note that the setting of $\lambda \ge 2L_f / \mu  $ implies $\Vert \nabla_{xy}^2 \fL(\,\cdot\,,\,\cdot\,) \Vert \le 2 \lambda L_g$, then we further have
\begin{align*}
    &\quad \Vert \nabla y_{\lambda}^*(x) - \nabla y^*(x) \Vert \\
    & \le\left  \Vert \nabla_{xy}^2 g(x,y^*(x)) - \frac{\nabla_{xy}^2 \fL_{\lambda}(x,y_{\lambda}^*(x))}{ \lambda} \right \Vert \left \Vert [\nabla_{yy}^2 g(x,y^*(x))]^{-1}  \right \Vert  \\
    &\quad + \left \Vert \frac{\nabla_{xy}^2 \fL_{\lambda}(x,y_{\lambda}^*(x))}{\lambda}  \right \Vert  \left \Vert [\nabla_{yy}^2 g(x,y^*(x))]^{-1} - \left [  \frac{\nabla_{yy}^2 \fL_{\lambda}(x,y_{\lambda}^*(x))}{\lambda} \right]^{-1}  \right \Vert  \\
    &\le \frac{1}{\mu} \left \Vert \nabla_{xy}^2 g(x,y^*(x)) - \nabla_{xy}^2 g(x,y_{\lambda}^*(x)) - \frac{\nabla_{xy}^2 f(x,y_{\lambda}^*(x))}{\lambda}  \right \Vert  +  \frac{2L_g}{\lambda \mu^2} \left(L_f + 
     \frac{C_f \rho_g}{\mu}
     \right)  \\
    &\le \left( \frac{1}{\lambda \mu} + \frac{2L_g}{\lambda \mu^2} \right) \left( L_f + \frac{C_f \rho_g}{\mu} \right).
    \end{align*}
\end{proof}

It is clear that $\Vert \nabla y^*(x) \Vert \le L_g/\mu$, therefore $y^*(x)$ is $(L_g/\mu)$-Lipschitz. Below, we show that a similar result also holds for $y_{\lambda}^*(x)$.

\begin{lem} \label{lem:nabla-y-lambda-bound}
Under Assumption \ref{asm:sc}, for $\lambda \ge 2L_f /\mu$, it holds that $\Vert \nabla y_{\lambda}^*(x) \Vert \le 4L_g/\mu$.    
\end{lem}


\begin{proof}
    Recall \Eqref{eq:y-lambda-x} that 
\begin{align*}
    \nabla y_{\lambda}^*(x) = -  \nabla_{xy}^2 \fL_{\lambda}(x,y_{\lambda}^*(x)) [\nabla_{yy}^2 \fL_{\lambda}(x,y_{\lambda}^*(x)) ]^{-1}.
\end{align*}
We arrive at the conclusion by noting 
that $  \nabla_{xy}^2 \fL_{\lambda}(\,\cdot\,,\,\cdot\,) \preceq 2 \lambda L_g$ and  $ \nabla_{yy}^2 \fL_{\lambda}(\,\cdot\, , \,\cdot\,) \succeq \lambda \mu /2 $ by Lemma \ref{lem:Lag-sc}.

\end{proof}
This implies that $y_{\lambda}^*(x)$ is $(4L_g/\mu)$-Lipschitz.

\begin{lem} \label{lem:nabla2-bound}
Under Assumption \ref{asm:sc}, for $\lambda \ge 2L_f /\mu$, $ \nabla \fL_{\lambda}^*(x)$ is $D_3$-Lipschitz, where
\begin{align*}
    D_3 &:= L_f + \frac{4L_f L_g}{\mu} +\frac{C_f \rho_g}{\mu} +  \frac{C_f L_g \rho_g}{\mu^2} + L_g D_2 = \fO(\kappa^3).
\end{align*}
\end{lem}

\begin{proof}
Note that
\begin{align} \label{eq:nabla-AB}
    \nabla \fL_{\lambda}^*(x) = \underbrace{\nabla_x f(x,y_{\lambda}^*(x))}_{A(x)} + \underbrace{\lambda ( \nabla_x g(x,y_{\lambda}^*(x)) - \nabla_x g(x,y^*(x)))}_{B(x)}. 
\end{align}
where $A(x)$ and $B(x)$ are both mappings  $\BR^{d_x} \rightarrow \BR^{d_x}$. 
From Lemma \ref{lem:nabla-y-lambda-bound} we know that $y_{\lambda}^*(x)$ is $(4L_g/\mu)$-Lipschitz. This further implies that $A(x)$ is  $(1+4L_g/\mu)L_f$-Lipschitz. Next, we bound the Lipschitz  coefficient of $B(x)$ via its derivative $\nabla B(x): \BR^{d_x} \rightarrow \BR^{d_x} \times \BR^{d_x} $, which has the following form by 
taking total derivative on $B(x)$:
\begin{align} \label{eq:nabla-Bx}
\begin{split}
    \nabla B(x) &= 
    \lambda (\nabla_{xx}^2 g(x,y_{\lambda}^*(x))-\nabla_{xx}^2 g(x,y^*(x))) \\
    &\quad + \lambda \left(\nabla y_{\lambda}^*(x)  \nabla_{yx}^2 g(x,y_{\lambda}^*(x)) - \nabla y^*(x) \nabla_{yx}^2 g(x,y^*(x)) \right).
\end{split}
\end{align}
And we can bound the operator norm of $\nabla B(x)$ by:
\begin{align*}
    \Vert \nabla B(x) \Vert &\le  \lambda \Vert \nabla_{xx}^2 g(x,y_{\lambda}^*(x)) - \nabla_{xx}^2 g(x,y^*(x)) \Vert \\
    &\quad + \lambda \Vert \nabla y^*(x) \Vert \Vert \nabla_{yx}^2 g(x,y_{\lambda}^*(x)) - \nabla_{yx}^2 g(x,y^*(x)) \Vert \\
    &\quad  + \lambda \Vert \nabla y_{\lambda}^*(x) - \nabla y^*(x) \Vert \Vert \nabla_{yx}^2 g(x,y_{\lambda}^*(x)) \Vert\\
    &\le  \lambda \rho_g \left(1+ \frac{L_g}{\mu} \right) \Vert y_{\lambda}^*(x) - y^*(x) \Vert +\lambda L_g \Vert \nabla y_{\lambda}^*(x) - \nabla y^*(x) \Vert  \\
    &\le  \left(1+ \frac{L_g}{\mu} \right) \frac{C_f \rho_g}{\mu} +L_g D_2,
\end{align*}
where we use Lemma \ref{lem:nabla-y-lambda-bound}
in the second inequality; Lemma \ref{lem:yy-lambda} and \ref{lem:yxyx-lambda} in the third one.
\end{proof}


\section{Lemmas for Finding Second-Order Stationarity} \label{apx:lem-2nd}


\begin{lem} \label{lem:nabla2-yx}
Under Assumption \ref{asm:sc} and \ref{asm:third}, 
for $\lambda \ge  2L_f / \mu$,
we have $ \Vert \nabla^2 y^*(x) -\nabla^2 y_{\lambda}^*(x)  \Vert \le D_4 /\lambda$, where 
\begin{align*}
    D_4 := \frac{2\rho_g}{\mu^2}  \left( 1 + \frac{L_g}{\mu}\right)^2  \left(L_f + 
     \frac{C_f \rho_g}{\mu}
     \right) +   \frac{14 L_g \rho_g D_2}{\mu^2} + \frac{50 L_g^2 }{\mu^3} \left( \frac{C_f \nu_g}{\mu} + \rho_f \right) = \fO(\kappa^5).
\end{align*}
\end{lem}

\begin{proof}
{First of all, we calculate the explicit form of $\nabla^2 y^*(x)$ and $\nabla^2 y_{\lambda}^*(x)$.}

    By taking the derivative with respect to $x$ on
    \begin{align*}
        \nabla_{xy}^2 g(x,y^*(x)) + \nabla y^*(x) \nabla_{yy}^2 g(x,y^*(x)) =0,
    \end{align*}
    we obtain
    \begin{align*}
        & \nabla_{xxy}^3 g(x,y^*(x)) + \nabla_{yxy}^3 g(x,y^*(x)) \times_{1}  \nabla y^*(x) + \nabla^2 y^*(x) \times_3 \nabla_{yy}^2 g(x,y^*(x)) \\
        &\quad + \nabla_{xyy}^3 g(x,y^*(x)) \times_2 \nabla y^*(x)  +\nabla_{yyy}^3 g(x,y^*(x)) \times_1 \nabla y^*(x) \times_2 \nabla y^*(x) = 0.
    \end{align*}
    Rearranging to get
    {\begin{align} \label{eq:nabla2-yx}
    \begin{split}
         \nabla^2 y^*(x) 
        &= -\left( \nabla_{xxy}^3 g(x,y^*(x)) + \nabla_{yxy}^3 g(x,y^*(x)) \times_1 \nabla y^*(x) 
        \right)  \times_3 [\nabla_{yy}^2 g(x,y^*(x))]^{-1} \\
        &\quad -
        \nabla_{xyy}^3 g(x,y^*(x)) \times_2 \nabla y^*(x) \times_3 [\nabla_{yy}^2 g(x,y^*(x))]^{-1} \\
        &\quad-  \nabla_{yyy}^3 g(x,y^*(x)) \times_1 \nabla y^*(x) \times_2 \nabla y^*(x)    \times_3 [\nabla_{yy}^2 g(x,y^*(x))]^{-1}.
    \end{split}
    \end{align}}
    Similarly,
    {
    \begin{align} \label{eq:nabla2-yx-lambda}
    \begin{split}
        \nabla^2 y_{\lambda}^*(x)  &= -\left( \nabla_{xxy}^3 \fL_{\lambda} (x, y_{\lambda}^*(x)) + \nabla_{yxy}^3 \fL_\lambda(x, y_{\lambda}^*(x)) \times_1 \nabla y_{\lambda}^*(x)  \right) \times_3
        [ \nabla_{yy}^2 \fL_{\lambda}(x,y_{\lambda}^*(x))  ]^{-1} \\
        &\quad - 
        \nabla_{xyy}^3 \fL_{\lambda}(x,y_{\lambda}^*(x))
        \times_2 \times_3 [ \nabla_{yy}^2 \fL_{\lambda}(x,y_{\lambda}^*(x))]^{-1} \\
        &\quad - 
        \nabla y_{\lambda}^*(x)   +  
         \nabla_{yyy}^3 \fL_{\lambda} (x,y_{\lambda}^*(x))
        \times_1 \nabla y_{\lambda}^*(x)
        \times_2 \nabla y_{\lambda}^*(x) 
         \times_3 [ \nabla_{yy}^2 \fL_{\lambda}(x,y_{\lambda}^*(x))]^{-1}.
    \end{split}
    \end{align}}
{We then prove $\Vert \nabla^2 y^*(x) - \nabla^2 y_{\lambda}^*(x) \Vert = \fO(1/\lambda)$. We prove this by showing that the difference between each corresponding term of $\nabla^2 y^*(x)$ and $\nabla^2 y_{\lambda}^*(x)$ is $\fO(1/\lambda)$.}
Note that 
\begin{align*}
    \left \Vert  \nabla_{xxy}^3  g(x,y^*(x)) - \frac{\nabla_{xxy}^3 \fL_{\lambda} (x,y_{\lambda}^*(x))  }{\lambda}   \right \Vert \le \nu_g \Vert y_{\lambda}^*(x) - y^*(x) \Vert + \frac{\rho_f}{\lambda} \le \frac{1}{\lambda} \left(\frac{C_f \nu_g}{\mu} +\rho_f \right),
\end{align*}
and
\begin{align*}
    &\quad \left \Vert 
    \nabla_{yxy}^3 g(x,y^*(x)) \times_1 \nabla y^*(x)  - \frac{\nabla_{yxy}^3 \fL_{\lambda} (x,y_{\lambda}^*(x)) \times_1 \nabla y_{\lambda}^*(x) }{\lambda}
    \right \Vert \\
    &\le \Vert \nabla y^*(x) - \nabla y_{\lambda}^*(x) \Vert \Vert \nabla_{yxy}^3 g(x,y^*(x)) \Vert + \Vert \nabla y_{\lambda}^*(x) \Vert \left\Vert \nabla_{yxy}^3 g(x,y^*(x)) - \frac{\nabla_{yxy}^3 \fL_{\lambda}(x, y_{\lambda}^*(x))}{\lambda} \right \Vert \\
    &\le \frac{\rho_g D_2}{\lambda} + \frac{4L_g}{\lambda \mu} \left( \frac{C_f \nu_g}{\mu} + \rho_f \right),
\end{align*}
and
{ 
\begin{align*}
    &\quad  \left \Vert
   \nabla_{yyy}^3 g(x,y^*(x)) \times_1  \nabla y^*(x)  \times_2 \nabla y^*(x)  - \frac{\nabla_{yyy}^3 \fL_{\lambda} (x,y_{\lambda}^*(x)) \times_1 \nabla y_{\lambda}^*(x) \times_2  \nabla y_{\lambda}^*(x)} { \lambda}
    \right \Vert \\
    &\le \Vert \nabla y^*(x) \Vert \Vert \nabla_{yyy}^3 g(x,y^*(x))\Vert \Vert \nabla y^*(x) - \nabla y_{\lambda}^*(x) \Vert  \\
    &\quad + \Vert \nabla y_{\lambda}^*(x) \Vert \Vert \nabla_{yyy}^3 g(x,y^*(x))\Vert \Vert \nabla y^*(x) - \nabla y_{\lambda}^*(x) \Vert  \\
    &\quad +  \Vert \nabla y_{\lambda}^*(x) \Vert^2  \left \Vert  \nabla_{xxy}^3  g(x,y^*(x)) - \frac{\nabla_{xxy}^3 \fL_{\lambda} (x,y_{\lambda}^*(x))  }{\lambda}   \right \Vert \\
    &\le  \frac{5L_g \rho_g D_2}{\lambda \mu} + \frac{16L_g^2 }{\lambda \mu^2} \left( \frac{C_f \nu_g}{\mu} + \rho_f \right),
\end{align*}}
we can obtain that
{
\begin{align*}
    &\quad \Vert \nabla^2 y^*(x) - \nabla^2 y_{\lambda}^*(x) \Vert  \\
    &\le \rho_g \left( 1 + \frac{L_g}{\mu}\right)^2 \left \Vert [\nabla_{yy}^2 g(x,y^*(x))]^{-1} - \left [  \frac{\nabla_{yy}^2 \fL_{\lambda}(x,y_{\lambda}^*(x))}{\lambda} \right]^{-1}  \right \Vert \\
    &\quad + \left( \frac{7L_g \rho_g D_2}{\lambda \mu} + \frac{25L_g^2 }{\lambda \mu^2} \left( \frac{C_f \nu_g}{\mu} + \rho_f \right) \right) \left \Vert \left [  \frac{\nabla_{yy}^2 \fL_{\lambda}(x,y_{\lambda}^*(x))}{\lambda} \right]^{-1} \right \Vert \\
    &\le \frac{2\rho_g}{\lambda \mu^2}  \left( 1 + \frac{L_g}{\mu}\right)^2  \left(L_f + 
     \frac{C_f \rho_g}{\mu}
     \right) +   \frac{14 L_g \rho_g D_2}{\lambda \mu^2} + \frac{50 L_g^2 }{\lambda \mu^3} \left( \frac{C_f \nu_g}{\mu} + \rho_f \right) ,
\end{align*}}
where we use \Eqref{eq:diff-inv} in the second inequality.
\end{proof}

\begin{lem} \label{lem:nabla3-bound}
Under Assumption \ref{asm:sc} and \ref{asm:third}, for $ \lambda \ge 2 L_f /\mu$, 
$\nabla^2 \fL_{\lambda}^*(x) $ is $D_5$-Lipschitz,
 where
\begin{align*}
    D_5 &:=  \left( 1 + \frac{4L_g}{\mu} \right)^2 \left( 3\rho_f  + \frac{2 L_f \rho_g}{\mu} \right) + \left(1 + \frac{L_g}{\mu} \right)^2 \frac{C_f \nu_g}{\mu} \\
    &\quad + \left(2+ \frac{5L_g}{\mu} \right) D_2 \rho_g +  \left( 1+ \frac{L_g}{\mu} \right)^2 \frac{C_f \rho_g^2}{\mu^2} + L_g D_4   = \fO(\ell \kappa^5).
\end{align*}
\end{lem}

\begin{proof}
Similar to the proof of Lemma \ref{lem:nabla2-bound},  we split $\nabla^2 \fL_{\lambda}^*(x)$ into two terms:
\begin{align*}
    \nabla^2 \fL_{\lambda}^*(x) = \nabla A(x) + \nabla B(x),
\end{align*}
where both the mappings $A(x)$ and $B(x)$ both follow the definitions in \Eqref{eq:nabla-AB}. 
rTaking total derivative on $A(x)$, we obtain
\begin{align*}
    \nabla A(x) = \nabla_{xx}^2 f(x,y_{\lambda}^*(x)) + \nabla y_{\lambda}^*(x) \nabla_{yx}^2 f(x,y_{\lambda}^*(x)).
\end{align*}
And recall \Eqref{eq:nabla-Bx} that 
\begin{align*}
    \nabla B(x) &= \lambda (\nabla_{xx}^2 g(x,y_{\lambda}^*(x))-\nabla_{xx}^2 g(x,y^*(x))) \\
    &\quad + \lambda \left(\nabla y_{\lambda}^*(x)  \nabla_{yx}^2 g(x,y_{\lambda}^*(x)) - \nabla y^*(x) \nabla_{yx}^2 g(x,y^*(x)) \right).
\end{align*}
Then we bound the Lipschitz coefficient of $\nabla A(x)$ and $\nabla B(x)$, respectively. 

From \Eqref{eq:nabla2-yx} we can calculate that 
\begin{align} \label{bound:nabla2-yx}
    \Vert \nabla^2 y^*(x) \Vert \le \left( 1 + \frac{L_g}{\mu} \right)^2 \frac{\rho_g}{\mu}  .
\end{align}
Similarly, from \Eqref{eq:nabla2-yx-lambda} we can calculate that
\begin{align} \label{bound:nabla2-y-lambda}
    \Vert \nabla^2 y_{\lambda}^*(x)  \Vert &\le \left( 1 + \frac{4L_g}{\mu} \right)^2 \left( \frac{\rho_f}{\lambda}+ \rho_g \right) \frac{2}{\mu} \le \left( 1 + \frac{4L_g}{\mu} \right)^2 \left( \frac{2\rho_f}{L_f} + \frac{2 \rho_g}{\mu} \right). 
\end{align}
From Lemma \ref{lem:nabla-y-lambda-bound} we know that $y_{\lambda}^*(x)$ is $(4L_g/\mu)$-Lipschitz. This further implies that both $ \nabla_{xx}^2 f(x,y_{\lambda}^*(x))$ and $\nabla_{yx}^2 f(x,y_{\lambda}^*(x)) $ are $(1+ 4L_g/ \mu)\rho_f$-Lipschitz. Then, for any $x_1,x_2 \in \BR^{d_x}$, we have
\begin{align*}
    &\quad  \Vert \nabla A(x_1) - \nabla A(x_2) \Vert \\
   &\le \Vert  \nabla_{xx}^2 f(x_1,y_{\lambda}^*(x_1)) - \nabla_{xx}^2 f(x_2, y_{\lambda}^*(x_2)) \Vert  \\
   &\quad + \Vert \nabla y_{\lambda}^*(x_1) \nabla_{yx}^2 f(x_1,y_{\lambda}^*(x_1)) - \nabla y_{\lambda}^*(x_2) \nabla_{yx}^2 f(x_1,y_{\lambda}^*(x_1)) \Vert \\
   &\quad + \Vert \nabla y_{\lambda}^*(x_2) \nabla_{yx}^2 f(x_1,y_{\lambda}^*(x_1)) - \nabla y_{\lambda}^*(x_2) \nabla_{yx}^2 f(x_2,y_{\lambda}^*(x_2)) \Vert \\
   &\le \Vert \nabla_{xx}^2 f(x_1,y_{\lambda}^*(x_1)) - \nabla_{xx}^2 f(x_2, y_{\lambda}^*(x_2)) \Vert  \\
   &\quad + \Vert \nabla y_{\lambda}^*(x_1) - \nabla y_{\lambda}^*(x_2) \Vert \Vert \nabla_{yx}^2 f(x_1,y_{\lambda}^*(x_1)) \Vert \\
   &\quad + \Vert \nabla y_{\lambda}^*(x_2) \Vert \Vert \nabla_{yx}^2 f(x_1,y_{\lambda}^*(x_1)) - \nabla_{yx}^2 f(x_2,y_{\lambda}^*(x_2)) \Vert \\
   &\le \underbrace{\left(\left( 1+ \frac{4L_g}{\mu} \right)^2 \rho_f + \left( 1 + \frac{4L_g}{\mu} \right)^2 \left( \frac{2\rho_f}{L_f} + \frac{2 \rho_g}{\mu} \right) L_f \right)}_{C_1}\Vert x_1-x_2 \Vert,
\end{align*}
where $C_1$ gives the upper bound of the Lipschitz coefficient of the mapping $\nabla A(x)$.

To bound the Lipschitz coefficient of $\nabla B(x)$,
we first derive the explicit form of the mapping $\nabla^2 B(x): \BR^{d_x} \rightarrow \BR^{d_x} \times \BR^{d_x} \times \BR^{d_x} $ by:
\begin{align*}
    \nabla^2 B(x) 
    &= 
    \nabla^2 y_{\lambda}^*(x) \times_3 [\nabla_{yx}^2 f(x,y_{\lambda}^*(x))]^\top \\
    &\quad + \lambda ( \nabla_{xxx}^3 g(x,y_{\lambda}^*(x)) - \nabla_{xxx}^3 g(x,y^*(x)) ) \\
    &\quad + \lambda ( 
    \nabla_{yxx}^3 g(x,y_{\lambda}^*(x)) \times_1 \nabla y_{\lambda}^*(x) - \nabla_{yxx}^3 g(x,y^*(x)) \times_1 \nabla y^*(x) 
    ) \\
    &\quad + \lambda \left(  \nabla_{xyx}^3 g(x,y_{\lambda}^*(x)) \times_2 \nabla y_{\lambda}^*(x) - \nabla_{xyx}^3 g(x,y^*(x))   \times_2 \nabla y^*(x) \right)  \\
    &\quad + \lambda \left( 
     \nabla_{yyx}^3 g(x,y_{\lambda}^*(x)) \times_1 \nabla y_{\lambda}^*(x) \times_2  \nabla y_{\lambda}^*(x)   - 
    \nabla_{yyx}^3 g(x,y^*(x)) \times_1 \nabla y^*(x)  \times_2 \nabla y^*(x)
    \right) \\
    &\quad + \lambda \left(  \nabla^2 y_{\lambda}^*(x) \times_3 [\nabla_{yx}^2 g(x,y_{\lambda}^*(x))]^\top -  \nabla^2 y^*(x) \times_3 [\nabla_{yx}^2 g(x,y^*(x))]^\top\right).
\end{align*}
Then we can bound the Lipschitz coefficient of $\nabla B(x)$
via the operator norm of $\nabla^2 B(x)$:  
\begin{align*}
   C_2&:= \Vert \nabla^2 B(x) \Vert \\
    &\le 
    \Vert \nabla_{xxx}^3 g(x,y^*(x)) - 
    \nabla_{xxx}^3 g(x,y_{\lambda}^*(x))\Vert  \\
    &\quad + \lambda \Vert \nabla y^*(x) \Vert \Vert \nabla_{yxx}^3 g(x,y^*(x)) - \nabla_{yxx}^3 g(x,y_{\lambda}^*(x)) \Vert \\
    &\quad+ \lambda \Vert \nabla y_{\lambda}^*(x) - \nabla y^*(x) \Vert \Vert \nabla_{yxx}^3 g(x,y_{\lambda}^*(x)) \Vert \\
    &\quad + \lambda \Vert \nabla y^*(x) \Vert \Vert \nabla_{xyx}^3 g(x,y^*(x)) - \nabla_{xyx}^3 g(x,y_{\lambda}^*(x)) \Vert  \\
    &\quad + \lambda \Vert \nabla y^*(x) - \nabla y_{\lambda}^*(x) \Vert \Vert \nabla_{xyx}^3 g(x,y_{\lambda}^*(x)) \Vert \\
    &\quad + \lambda \Vert \nabla y^*(x) \Vert \Vert \nabla_{yyx}^3 g(x,y^*(x)) \Vert  \Vert \nabla y_{\lambda}^*(x) - \nabla y^*(x) \Vert  \\
    &\quad + \lambda \Vert \nabla y_{\lambda}^*(x) \Vert \Vert \nabla_{yyx}^3 g(x,y^*(x)) \Vert  \Vert \nabla y_{\lambda}^*(x) - \nabla y^*(x) \Vert \\
    &\quad + \lambda \Vert \nabla y^*(x) \Vert^2 \Vert \nabla_{yyx}^3 g(x,y^*(x)) - \nabla_{yyx}^3 g(x,y_{\lambda}^*(x)) \Vert \\
    &\quad + \lambda \Vert \nabla^2 y^*(x) \Vert \Vert \nabla_{yx}^2 g(x,y^*(x)) - \nabla_{yx}^2 g(x,y_{\lambda}^*(x)) \Vert \\
    &\quad + \lambda \Vert \nabla^2 y^*(x) - \nabla^2 y_{\lambda}^*(x) \Vert  \Vert \nabla_{yx}^2 g(x,y_{\lambda}^*(x)) \Vert,
\end{align*}
{which only requires using triangle inequality multiple times. Then we plug
\begin{align} \label{eq:lambda}
\begin{split} 
            \Vert y_{\lambda}^*(x) - y^*(x) \Vert &= \fO(1/\lambda), \text{ By Lemma \ref{lem:yy-lambda} } \\
        \Vert \nabla y_{\lambda}^*(x) - \nabla y^*(x) \Vert &= \fO(1/\lambda), \text{ By Lemma \ref{lem:yxyx-lambda}. } \\
        \Vert \nabla^2 y_{\lambda}^*(x) - \nabla^2 y^*(x) \Vert &= \fO(1/\lambda), \text{ By Lemma \ref{lem:nabla2-yx}.}
\end{split}
\end{align}
and
\begin{align} \label{eq:1}
\begin{split}
        \Vert \nabla y^*(x) \Vert &= \fO(1) , \text{ By \Eqref{eq:y-x}.} \\
    \Vert \nabla y_{\lambda}^*(x) \Vert &= \fO(1), \text{ By Lemma \ref{lem:nabla-y-lambda-bound}}. \\
    \Vert \nabla^2 y^*(x) \Vert &= \fO(1), \text{ By \Eqref{bound:nabla2-yx}}.
\end{split}
\end{align}
into the bound for $C_2$ to obtain that
} 
\begin{align*}
    C_1+C_2 & \le \left( 1 + \frac{4L_g}{\mu} \right)^2 \left( 3\rho_f  + \frac{2 L_f \rho_g}{\mu} \right) + \left(1 + \frac{L_g}{\mu} \right)^2 \frac{C_f \nu_g}{\mu} \\
    &\quad + \left(2+ \frac{5L_g}{\mu} \right) D_2 \rho_g +  \left( 1+ \frac{L_g}{\mu} \right)^2 \frac{C_f \rho_g^2}{\mu^2} + L_g D_4.
\end{align*}
\end{proof}
 
\begin{lem} \label{lem:nabla2-phi}
Under Assumption \ref{asm:sc} and \ref{asm:third}, for $\lambda \ge 2 L_f /\mu$, we have $ \Vert \nabla^2 \fL_{\lambda}^*(x) - \nabla^2 \varphi(x) \Vert \le  D_6 /\lambda$,
where
\begin{align*}
    D_6 := 2L_g D_2^2 + \left( 1 + \frac{L_g}{\mu}\right)^2\left( \frac{C_f \rho_f}{\mu} + \frac{C_f L_f \rho_g}{\mu^2} + \frac{C_f^2 \nu_g}{2 \mu^2} + \frac{C_f^2 \rho_g^2}{2 \mu^3}  \right)  = \fO(\ell \kappa^6).
\end{align*}
\end{lem}

\begin{proof}
    Taking total derivative on $ \nabla \varphi(x) = \nabla_x f(x,y^*(x)) + \nabla y^*(x) \nabla_y f(x,y^*(x))$ yields
    \begin{align} \label{eq:phi-hess}
\begin{split}
    \nabla^2 \varphi(x) &= \nabla_{xx}^2 f(x,y^*(x)) + \nabla y^*(x) \nabla_{yx}^2 f(x,y^*(x))  + \nabla^2 y^*(x) \bar \times_3 \nabla_y f(x,y^*(x)) \\
    &\quad + \nabla_{xy}^2f(x,y^*(x)) [\nabla y^*(x)]^\top + \nabla y^*(x) \nabla_{yy}^2 f(x,y^*(x)) [\nabla y^*(x)]^\top.
\end{split}
\end{align}
Plug into the close form of $\nabla^2 y^*(x)$ given by \Eqref{eq:nabla2-yx}, we arrive at
    \begin{align*}
    &\quad \nabla^2 \varphi(x) \\
    &= \underbrace{\nabla_{xx }^2  f(x,y^*(x)) - \nabla_{xxy}^3 g(x,y^*(x)) \times_3 [  \nabla_{yy}^2 g(x,y^*(x))]^{-1} \bar \times_3 \nabla_y f(x,y^*(x)) }_{\rm (I)} \\
    &\quad +  
    \underbrace{  \left(\nabla_{yx}^2 f(x,y^*(x)) - \nabla_{yxy}^3 g(x,y^*(x)) \times_3 [\nabla_{yy}^2 g(x,y^*(x)]^{-1}  \bar \times_3 \nabla_y f(x,y^*(x)) \right) \times_1 \nabla y^*(x)  }_{\rm (II)}  \\
    &\quad +  
    \underbrace{  \left(\nabla_{xy}^2 f(x,y^*(x)) - \nabla_{xyy}^3 g(x,y^*(x)) \times_3 [\nabla_{yy}^2 g(x,y^*(x)]^{-1}  \bar \times_3 \nabla_y f(x,y^*(x)) \right) \times_2 \nabla y^*(x)  }_{\rm (III)}  \\
    &\quad + \underbrace{  \left( \nabla_{yy}^2 f(x,y^*(x)) -  \nabla_{yyy}^3 g(x,y^*(x)) \times_3 [\nabla_{yy}^2 g(x,y^*(x)]^{-1} \bar \times_3 \nabla_y f(x,y^*(x))   \right) \times_1 \nabla y^*(x) \times_2 \nabla y^*(x) }_{\rm (IV)}. 
\end{align*}
Recall that
\begin{align*}
    \nabla^2 \fL_{\lambda}^*(x) &= \nabla_{xx}^2  f(x,y_{\lambda}^*(x)) +  \lambda ( \nabla_{xx}^2 g(x,y_{\lambda}^*(x)) - \nabla_{xx}^2 g(x,y^*(x))) \\
    &\quad + \nabla y_{\lambda}^*(x) \nabla_{yx}^2 \fL_{\lambda}(x,y_{\lambda}^*(x))
    - \lambda \nabla y^*(x) \nabla_{yx}^2 g(x,y^*(x)) .
\end{align*}
{Our goal is show that $\nabla \fL_{\lambda}^*(x) \approx \nabla^2 \varphi(x)$. At first glance, these two quantities are very different and we can not directly bound their difference: $\nabla^2 \varphi(x)$ takes the form of $A + B C + C^\top B^\top + B D B^\top $, while $\nabla^2 \fL_{\lambda}^*(x)$ looks different.  Below, we introduce an
intermediary quantity $\tilde \nabla^2 \fL_{\lambda}^*(x)$ which takes a similar form as $\nabla^2 \varphi(x)$ to serves as a bridge:
}
\begin{align} \label{eq:inter-term}
\begin{split}
        \tilde \nabla^2 \fL_{\lambda}^*(x) &= \underbrace{\nabla_{xx}^2  f(x,y_{\lambda}^*(x)) +  \lambda ( \nabla_{xx}^2 g(x,y_{\lambda}^*(x)) - \nabla_{xx}^2 g(x,y^*(x)))}_{\rm (I')}\\
    &+  \underbrace{\nabla y^*(x) \left(\nabla_{yx}^2 \fL_{\lambda}(x,y_{\lambda}^*(x))- \lambda \nabla_{yx}^2 g(x,y^*(x))\right) }_{\rm (II')}\\ 
    & + \underbrace{ \left(\nabla_{xy}^2 \fL_{\lambda}(x,y_{\lambda}^*(x))- \lambda \nabla_{xy}^2 g(x,y^*(x))\right)[\nabla y^*(x)]^\top }_{\rm (III')} \\
    &+ \underbrace{\nabla y^*(x) \left(\nabla_{yy}^2 \fL_{\lambda} (x,y_{\lambda}^*(x))   - \lambda \nabla_{yy}^2 g(x,y^*(x)) \right) [\nabla y^*(x)]^\top }_{\rm (IV')}. 
\end{split}
\end{align}
{Now we show $\nabla^2 \fL_{\lambda}^*(x) \approx \tilde \nabla^2 \fL_{\lambda}^*(x) \approx \nabla^2 \varphi(x)$.}
% Recall  that
% \begin{align*}
%     \nabla y_{\lambda}^*(x) = - \nabla_{xy}^2 \fL_{\lambda}(x,y_{\lambda}^*(x)) [\nabla_{yy}^2 \fL_{\Ulambda}(x,y_{\lambda}^*(x)) ]^{-1}.
% \end{align*}
\begin{align*}
    &\quad \tilde \nabla^2 \fL_{\lambda}^*(x) - \nabla^2 \fL_{\lambda}^*(x) \\
    &= ( \nabla y^*(x) - \nabla y_{\lambda}^*(x) ) \nabla_{yx}^2 \fL_{\lambda} (x,y_{\lambda}^*(x)) + \\
    &\quad + \nabla_{xy}^2 \fL_{\lambda}(x,y_{\lambda}^*(x)) [ \nabla y^*(x)]^\top + \nabla y^*(x) \nabla_{yy}^2 \fL_{\lambda} (x,y_{\lambda}^*(x)) [\nabla y^*(x)]^\top \\
    &= ( \nabla y^*(x) - \nabla y_{\lambda}^*(x) ) \nabla_{yx}^2 \fL_{\lambda} (x,y_{\lambda}^*(x)) + \nabla_{xy}^2 \fL_{\lambda}(x,y_{\lambda}^*(x)) [ \nabla y^*(x) - \nabla y_{\lambda}^*(x)]^\top \\
    &\quad + \nabla y^*(x) \nabla_{yy}^2 \fL_{\lambda} (x,y_{\lambda}^*(x)) [\nabla y^*(x)]^\top - \nabla y_{\lambda}^*(x) \nabla_{yy}^2 \fL_{\lambda} (x,y_{\lambda}^*(x)) [\nabla y_{\lambda}^*(x)]^\top \\
    &= ( \nabla y^*(x) - \nabla y_{\lambda}^*(x) ) \nabla_{yx}^2 \fL_{\lambda} (x,y_{\lambda}^*(x)) + \nabla_{xy}^2 \fL_{\lambda}(x,y_{\lambda}^*(x)) [ \nabla y^*(x) - \nabla y_{\lambda}^*(x)]^\top \\ 
    &\quad + (\nabla y^*(x) - \nabla y_{\lambda}^*(x)) \nabla_{yy}^2 \fL_{\lambda} (x,y_{\lambda}^*(x)) [ \nabla y_{\lambda}^*(x)]^\top \\
    &\quad + \nabla y_{\lambda}^*(x) \nabla_{yy}^2 \fL_{\lambda} (x,y_{\lambda}^*(x)) (\nabla y^*(x) - \nabla y_{\lambda}^*(x))^\top \\
    &\quad + (\nabla y^*(x) - \nabla y_{\lambda}^*(x)) \nabla_{yy}^2 \fL_{\lambda} (x,y_{\lambda}^*(x)) [ \nabla y^*(x) - \nabla y_{\lambda}^*(x)]^\top \\
    &= (\nabla y^*(x) - \nabla y_{\lambda}^*(x)) \nabla_{yy}^2 \fL_{\lambda} (x,y_{\lambda}^*(x)) [ \nabla y^*(x) - \nabla y_{\lambda}^*(x)]^\top,
\end{align*}
{where we use the identity
\begin{align*}
    U A U^\top - V A V^\top = (U-V) A V^\top + VA (U-V)^\top + (U-V)A (U-V)^\top
\end{align*}
in the second last step, and we cancel the first four terms by
\Eqref{eq:y-lambda-x} in the final step.} Noting that $ \nabla_{yy}^2 \fL_{\lambda}(\,\cdot\, , \,\cdot\,) \preceq 2 \lambda L_g$, we have
\begin{align*}
    \Vert \tilde \nabla^2 \fL_{\lambda}^*(x) - \nabla^2 \fL_{\lambda}^*(x) \Vert \le 2 \lambda L_g \Vert \nabla y^*(x) - \nabla y_{\lambda}^*(x) \Vert^2 \le \frac{2L_g D_2^2}{\lambda }. 
\end{align*}
{Now, we have successfully simplified our goal to showing $\tilde \nabla^2 \fL_{\lambda}^*(x) \approx \nabla^2 \varphi(x)$, which can be done by showing 
${\rm (I)} \approx {\rm (I')}$, ${\rm (II)} \approx {\rm (II')}$, ${\rm (III)} \approx {\rm (III')}$, and ${\rm (IV)} \approx {\rm (IV')}$, separately. }Firstly,
{\begin{align*}
    &\quad \rm (I) - {\rm (I')} \\
    &= \nabla_{xx}^2 f(x,y^*(x)) - \nabla_{xx}^2 f(x,y_{\lambda}^*(x))  \\
    & \quad + \nabla_{xxy}^3 g(x,y^*(x)) \times_3 [\nabla_{yy}^2 g(x,y^*(x))]^{-1} \bar \times_3 ( \nabla_y f(x,y_{\lambda}^*(x)) - \nabla_y f(x,y^*(x))) \\
    &\quad + \lambda \left( \nabla_{xx}^2 g(x,y^*(x)) -  \nabla_{xx}^2 g(x,y_{\lambda}^*(x))  + \nabla_{xxy}^3  g(x,y^*(x)) \bar \times_3
    (y_{\lambda}^*(x) - y^*(x)) \right) \\
    &\quad + \lambda \nabla_{xxy}^3 g(x,y^*(x)) \times_3 [\nabla_{yy}^2 g(x,y^*(x))]^{-1} 
 \\
 &~~~~~~~~~~~~~~~~~~~~~~~~~~~\bar \times_3 \left(\nabla_y g(x,y_{\lambda}^*(x)) 
    - \nabla_y g(x,y^*(x)) - \nabla_{yy}^2 g(x,y^*(x)) (y_{\lambda}^*(x) -y^*(x))  \right).
\end{align*}}
Therefore, 
\begin{align*}
    \Vert {\rm (I) } - {\rm (I')} \Vert &\le \rho_f \Vert y^*(x) - y_{\lambda}^*(x) \Vert  + \frac{L_f \rho_g}{\mu} \Vert y^*(x) - y_{\lambda}^*(x) \Vert \\
    &\quad + \frac{\lambda \nu_g}{2} \Vert y^*(x) - y_{\lambda}^*(x) \Vert^2 + \frac{\lambda \rho_g^2}{2\mu} \Vert y^*(x) - y_{\lambda}^*(x) \Vert^2  \\
    &\le \frac{C_f \rho_f}{\lambda \mu} + \frac{C_f L_f \rho_g}{\lambda \mu^2} +  \frac{C_f^2 \nu_g}{2\lambda \mu^2}+ \frac{C_f^2 \rho_g^2}{2 \lambda \mu^3} .
\end{align*}
Using $\Vert \nabla y^*(x) \Vert \le  L_g / \mu$, we can similarly bound the difference between ${\rm (II)}$ and ${\rm (II')}$ by
{ \begin{align*}
    &\quad \lVert{\rm (II)} - {\rm (II')}\rVert \\
    &\le 
      \Vert \nabla y^*(x) \Vert \Big \Vert  \nabla_{yx}^2 f(x,y^*(x)) - \nabla_{yxy}^3 g(x,y^*(x)) \times_3 [\nabla_{yy}^2 g(x,y^*(x)]^{-1} \bar \times_3 \nabla_y f(x,y^*(x))\\
     &\quad - \nabla_{yx}^2 f(x,y_{\lambda}^*(x))-\lambda  \left(\nabla_{yx}^2 g(x,y_{\lambda}^*(x))-\nabla_{yx}^2 g(x,y^*(x)) \right) \Big \Vert \\
     & \le \frac{L_g }{\mu} \left( \frac{C_f \rho_f}{\lambda \mu} + \frac{C_f L_f \rho_g}{\lambda \mu^2}  + \frac{C_f^2 \nu_g}{2\lambda \mu^2} +\frac{C_f^2 \rho_g^2}{2 \lambda \mu^3} \right),
\end{align*}}
And the bound for ${\rm(III)}$ and $ {\rm (III')}$ is the same. Finally, we know that
\begin{align*}
    &\quad \left \Vert {\rm (IV)} - {\rm (IV')} \right \Vert \\
    &\le \Vert \nabla y^*(x) \Vert^2 \Big \Vert \nabla_{yy}^2 f(x,y^*(x)) - \nabla_{yyy}^3 g(x,y^*(x)) \times_3 [ \nabla_{yy}^2 g(x,y^*(x))]^{-1} \bar \times_3 \nabla_y f(x,y^*(x))  \\
    &\quad - \nabla_{yy}^2 f(x,y_{\lambda}^*(x)) - \lambda \left( \nabla_{yy}^2 g(x,y_{\lambda}^*(x)) - \nabla_{yy}^2 g(x,y^*(x)) \right) \Big \Vert \\
    &\le \frac{L_g^2}{\mu^2} \left( \frac{C_f \rho_f}{\lambda \mu} + \frac{C_f L_f \rho_g}{\lambda \mu^2}  + \frac{C_f^2 \nu_g}{2\lambda \mu^2} +  \frac{C_f^2 \rho_g^2}{2 \lambda \mu^3}\right).
\end{align*}
Combing the above bounds completes our proof.
\end{proof}

\section{Proofs of Finding First-Order Stationarity} \label{apx:1st}

We recall the standard result for gradient descent for strongly convex functions, which is used for proving the linear convergence of $y_t^k \rightarrow y_{\lambda}^*(x_t) $ and $z_{t}^k \rightarrow y^*(x_t)$.
\begin{thm}[{\citet[Theorem 3.10]{bubeck2015convex}}] \label{thm:GD-SC}
Suppose $h(x): \BR^d \rightarrow \BR$ is  $\beta$-gradient Lipschitz and $\alpha$-strongly convex. Consider the following update of gradient descent:
\begin{align*}
    x_{t+1} = x_t - \frac{1}{\beta} \nabla h(x_t).
\end{align*}
Let $x^* = \arg \min_{x \in \BR^d} h(x)$. Then it holds that
\begin{align*}
    \Vert x_{t+1} - x^* \Vert^2 \le \left(1 - \frac{\alpha}{\beta} \right) \Vert x_t - x^* \Vert^2.
\end{align*}
\end{thm}

%\propGD*

% \begin{proof}
% It is a straightforward result of the following descent lemma for gradient descent.
% \begin{align*}
%     \fL_{\lambda}^*(x_{t+1}) &\le \fL_{\lambda}^*(x_t) + \langle \nabla \fL_{\lambda}^*(x_t) , x_{t+1} - x_t
%     \rangle + \frac{L}{2} \Vert x_{t+1 } - x_t \Vert^2 \\
%     &= \fL_{\lambda}^*(x_t) -
%     \frac{\eta}{2} \Vert \nabla \fL_{\lambda}^*(x_t) \Vert^2 -
%     \left( \frac{\eta}{2} - \frac{\eta^2 L}{2} \right) \Vert \hat \nabla \fL_{\lambda}^*(x_t) \Vert^2 + \frac{\eta}{2} \Vert \hat  \nabla \fL_{\lambda}^*(x_t) - \nabla \fL_{\lambda}^*(x_t) \Vert^2 \\
%     &\le \fL_{\lambda}^*(x_t) 
%     - \frac{\eta}{2} \Vert \nabla \fL_{\lambda}^*(x_t) \Vert^2 
%     + \frac{\eta}{2} \Vert  \hat \nabla \fL_{\lambda}^*(x_t) - \nabla \fL_{\lambda}^*(x_t) \Vert^2.
% \end{align*}
% A simple telescoping completes the proof.
% \end{proof}

% \propPGD*

% \begin{proof}
%     It is the same as the proof of Theorem 2.7 in \cite{huang2022efficiently}.
% \end{proof}



% \properror*

% \begin{proof}
% According to
% \begin{align*}
%     \Vert \hat \nabla \fL_{\lambda}^*(x_t) - \nabla \fL_{\lambda}^*(x_t) \Vert \le (L_f+ \lambda L_g) \Vert y_{t+1}- y_{\lambda}^*(x_t) \Vert + \lambda L_g \Vert z_{t+1} - y^*(x_t) \Vert,
% \end{align*}
% it suffices to let 
% \begin{align*}
% \Vert y_{t+1} - y_{\lambda}^*(x_t) \Vert & \le \frac{\zeta}{4\lambda L_g}   ~{\rm and}~
% \Vert z_{t+1} - y^*(x_t) \Vert \le \frac{\zeta}{2 \lambda L_g}
% \end{align*}
% holds for all $t$. Below, we only show how to control $\Vert y_{t+1} - y_{\lambda}^*(x_t) \Vert$ since the analysis for $ \Vert z_{t+1} - y^*(x_t) \Vert$ is very similar. 
% Since $\fL_{\lambda}(x,y)$ is   $(\lambda \mu /2)$-strongly convex and $(2 \lambda L_g)$-smooth in $y$.
% From Theorem 2.1.14 in \cite{nesterov2018lectures}, we know that
% \begin{align*}
%     \Vert y_{t+1} - y_{\lambda}^*(x_t) \Vert \le \left( 1 - \frac{\mu}{4L_g} \right)^{K} \Vert y_t - y^*(x_t) \Vert.
% \end{align*}
% This indicates that we can  set
% \begin{align*}
%     K_t \ge \frac{4 L_g}{\mu} \log \left( 
%     \frac{4 \lambda L_g \Vert y_t - y_{\lambda}^*(x_t) \Vert}{\zeta}
%     \right), \quad \forall 0 \le t \le T-1.
% \end{align*}
% Next, we need to bound $\Vert y_t - y_{\lambda}^*(x_t) \Vert$ for all $t$. For the case of $t=0$, we have
% \begin{align*}
%     \Vert y_0 - y_{\lambda}^*(x_0) \Vert \le \Vert y_0 - y^*(x_0) \Vert + \Vert y^*(x_0) - y_{\lambda}^*(x_0) \Vert \le \Vert y_0 - y^*(x_0) \Vert + \frac{C_f}{\lambda \mu}.
% \end{align*}
% For the case of $t \ge 1$, we can prove it by induction. Suppose that we have $\Vert y_{t} - y_{\lambda}^*(x_{t-1})\Vert \le \zeta / (4 \lambda L_g)$, then
% \begin{align*}
%     \Vert y_t - y_{\lambda}^*(x_t) \Vert  &\le  \Vert y_{t} - y_{\lambda}^*(x_{t-1}) \Vert + \Vert y_{\lambda}^*(x_{t}) - y_{\lambda}^*(x_{t-1}) \Vert \\
%     &\le \frac{\zeta}{4 \lambda L_g} + \frac{4 L_g}{\mu} \Vert x_t - x_{t-1} \Vert \\
%     &\le \frac{\zeta}{4 \lambda L_g} + \frac{4 L_g \eta}{\mu} \Vert \hat \nabla \fL_{\lambda}^*(x_{t-1})\Vert \\
%     &\le \frac{\zeta}{4 \lambda L_g} + \frac{4 L_g \eta \zeta}{\mu} + \frac{4 L_g \eta}{\mu} \Vert \nabla \fL_{\lambda}^*(x_{t-1})\Vert,
% \end{align*}
% where the second inequality uses the fact that $y_{\lambda}^*(\,\cdot\,)$ is $(4 L_g/\mu)$-Lipschitz.
%     Recall that
%     \begin{align*}
%         \nabla \fL_{\lambda}^*(x) = f(x,y_{\lambda}^*(x)) + \lambda ( 
%         g(x,y_{\lambda}^*(x) ) - g(x,y^*(x))
%         ).
%     \end{align*}
%     According to Lemma \ref{lem:yy-lambda}, we have
%     \begin{align*}
%         \Vert \nabla \fL_{\lambda}^*(x) \Vert \le C_f + \lambda L_g \Vert y_{\lambda}^*(x) - y^*(x) \Vert \le \left( 1 + \frac{L_g}{\mu}\right) C_f.
%     \end{align*}
% Thus,
% \begin{align*}
%     \Vert y_t - y_{\lambda}^*(x_t) \Vert \le \frac{\zeta}{4 \lambda L_g} + \frac{4 L_g \eta }{\mu} \left( \zeta +   \left( 1 + \frac{L_g}{\mu}\right) C_f\right).
% \end{align*}

% \end{proof}

Below, we prove that F${}^2$BA achieves the near-optimal first-order oracle complexity.

\thmFBA* 
\begin{proof}
Let $L$ be the gradient Lipschitz coefficient of $\fL_{\lambda}^*(x)$.
According to Lemma \ref{lem:1st} and the setting of $\lambda$, we have
\begin{enumerate}[label=\alph*.]
    \item $\sup_{x \in \BR^{d_x}} \Vert \nabla \fL_{\lambda}^*(x) - \nabla \varphi(x) \Vert = \fO(\epsilon)$. 
    \item $\fL_{\lambda}^*(x_0) - \inf_{x \in \BR^{d_x}}  \fL_{\lambda}^*(x) = \fO(\Delta)$.
    \item $ L:=\sup_{x \in \BR^{d_x}}\Vert \nabla^2 \fL_{\lambda}^*(x) \Vert = \fO(\ell \kappa^3)$.
\end{enumerate}
Due to Lemma \ref{lem:yy-lambda}, we also have $ \Vert y_0 - y_{\lambda}^*(x_0) \Vert^2 + \Vert y_0 - y^*(x_0) \Vert^2 =\fO(R)$.
Now it suffices to show that the algorithm can find an $\epsilon$-first-order stationary of $\fL_{\lambda}^*(x)$. 
We begin from the following descent lemma for gradient descent. Let $\eta_x \le 1/(2L)$, then
\begin{align} \label{eq:des}
\begin{split}
     \fL_{\lambda}^*(x_{t+1})
    &\le \fL_{\lambda}^*(x_t) + \langle \nabla \fL_{\lambda}^*(x_t) , x_{t+1} - x_t
    \rangle + \frac{L}{2} \Vert x_{t+1 } - x_t \Vert^2 \\
    &= \fL_{\lambda}^*(x_t) -
    \frac{\eta_x}{2} \Vert \nabla \fL_{\lambda}^*(x_t) \Vert^2 -
    \left( \frac{\eta_x}{2} - \frac{\eta_x^2 L}{2} \right) \Vert \hat \nabla \fL_{\lambda}^*(x_t) \Vert^2 + \frac{\eta_x}{2} \Vert \hat  \nabla \fL_{\lambda}^*(x_t) - \nabla \fL_{\lambda}^*(x_t) \Vert^2 \\
    &\le \fL_{\lambda}^*(x_t) 
    - \frac{\eta_x}{2} \Vert \nabla \fL_{\lambda}^*(x_t) \Vert^2 - 
    \frac{1}{4\eta_x} \Vert x_{t+1} - x_{t} \Vert^2
    + \frac{\eta_x}{2} \Vert  \hat \nabla \fL_{\lambda}^*(x_t) - \nabla \fL_{\lambda}^*(x_t) \Vert^2.
\end{split}
\end{align}
Note that
\begin{align*}
    \Vert \hat \nabla \fL_{\lambda}^*(x_t) - \nabla \fL_{\lambda}^*(x_t) \Vert \le 2 \lambda L_g \Vert y_{t}^K- y_{\lambda}^*(x_t) \Vert + \lambda L_g \Vert z_{t}^K - y^*(x_t) \Vert.
\end{align*}
By Theorem \ref{thm:GD-SC}, we have
\begin{align*}
    \Vert y_t^K - y_{\lambda}^*(x_t) \Vert^2 &\le \exp\left( - \frac{\mu K}{4 L_g} \right) \Vert y_t^0 - y_{\lambda}^*(x_t) \Vert^2 \\
     \Vert z_t^K - y^*(x_t) \Vert^2 &\le \exp\left( - \frac{\mu K}{L_g} \right) \Vert z_t^0 - y^*(x_t) \Vert^2
\end{align*}
Therefore, we have
\begin{align} \label{eq:err}
    \Vert \hat \nabla \fL_{\lambda}^*(x_t)- \nabla \fL_{\lambda}^*(x_t) \Vert^2 \le 4 \lambda^2 L_g^2 \exp\left( - \frac{\mu K}{4 L_g} \right) \left( \Vert y_t^0 - y_{\lambda}^*(x_t) \Vert^2 + \Vert z_t^0 - y^*(x_t) \Vert^2 \right)
\end{align}
By Young's inequality, 
\begin{align*}
    \Vert y_{t+1}^0 - y_{\lambda}^*(x_{t+1}) \Vert^2 & \le 2 \Vert y_t^K - y_{\lambda}^*(x_{t}) \Vert^2 + 2 \Vert y_{\lambda}^*(x_{t+1}) - y_{\lambda}^*(x_t) \Vert^2  \\
    &\le 2 \exp \left( - \frac{\mu K}{4 L_g} \right) \Vert y_t^0 - y_{\lambda}^*(x_t) \Vert^2 + \frac{32 L_g^2}{\mu^2} \Vert x_{t+1} - x_t \Vert^2,
\end{align*}
where the second inequality follows Lemma \ref{lem:nabla-y-lambda-bound} that $y_{\lambda}^*(x)$ is $(4 L_g/\mu)$-Lipschitz.  Similarly, we can derive the recursion about $\Vert z_t^0 - y^*(x_t) \Vert^2$. 

Put them together and let $K \ge 8 L_g/ \mu$, we have
\begin{align} \label{eq:delta_t}
    &\quad \Vert y_{t+1}^0 - y_{\lambda}^*(x_{t+1}) \Vert^2  + \Vert z_{t+1}^0 - y^*(x_{t+1}) \Vert^2 \\
    &\le 2 \exp \left( - \frac{\mu K}{4 L_g} \right) \left(\Vert y_t^0 - y_{\lambda}^*(x_t) \Vert^2 + \Vert z_t^0 - y^*(x_t) \Vert^2 \right) +\frac{34 L_g^2}{\mu^2} \Vert x_{t+1} - x_t \Vert^2 \\
    &\le \frac{1}{2} \left(\Vert y_t^0 - y_{\lambda}^*(x_t) \Vert^2 + \Vert z_t^0 - y^*(x_t) \Vert^2 \right) +\frac{34 L_g^2}{\mu^2} \Vert x_{t+1} - x_t \Vert^2.
\end{align}
Telescoping over $t$ yields
\begin{align*}
    &\quad \Vert y_t^0 - y_{\lambda}^*(x_t) \Vert^2 + \Vert z_t^0 - y^*(x_t) \Vert^2 \\
    &\le \underbrace{\left( \frac{1}{2} \right)^t \left(\Vert y_0 - y_{\lambda}^*(x_0) \Vert^2 + \Vert y_0 - y^*(x_0) \Vert^2 \right) + \frac{34 L_g^2}{\mu^2} \sum_{j=0}^{t-1}  \left( \frac{1}{2}\right)^{t-1-j} \Vert x_{j+1} - x_j \Vert^2}_{:=(*)}.
\end{align*}
Plug into \Eqref{eq:err}, which, in conjunction with \Eqref{eq:des}, yields that
\begin{align*}
     \fL_{\lambda}^*(x_{t+1}) &\le \fL_{\lambda}^*(x_t) 
    - \frac{\eta_x}{2} \Vert \nabla \fL_{\lambda}^*(x_t) \Vert^2 - 
    \frac{1}{4\eta_x} \Vert x_{t+1} - x_{t} \Vert^2 + 2 \eta_x \times \underbrace{ \lambda^2 L_g^2 \exp \left( - \frac{\mu K}{4L_g}\right)}_{:=\gamma} \times (*).
\end{align*}
Telescoping over $t$ further yields
\begin{align} \label{eq:final}
\begin{split}
    \frac{\eta_x}{2}\sum_{t=0}^{T-1} \Vert \nabla \fL_{\lambda}^*(x_t) \Vert^2 &\le \fL_{\lambda}^*(x_0) - \inf_{x \in \BR^{d_x}} \fL_{\lambda}^*(x) + 4 \eta_x \gamma \left(\Vert y_0 - y_{\lambda}^*(x_0) \Vert^2 + \Vert y_0 - y^*(x_0) \Vert^2 \right) \\
    &\quad - \left( \frac{1}{4 \eta_x} - \frac{136 \eta_x \gamma L_g^2}{\mu^2} \right) \sum_{t=0}^{T-1} \Vert x_{t+1} - x_t \Vert^2
\end{split}
\end{align}
Let $K = \fO(\kappa \log(\lambda \kappa)) = \fO(\kappa \log(\nicefrac{\ell \kappa}{\epsilon}))$ such that $\gamma$ is sufficiently small with
\begin{align*}
    \gamma \le \min \left\{ \frac{\mu^2}{1088 \eta_x^2 L_g^2}, \frac{1}{4 \eta_x}\right\}.
\end{align*}
Then we have,
\begin{align*}
    \frac{1}{T} \sum_{t=0}^{T-1} \Vert \nabla \fL_{\lambda}^*(x_t) \Vert^2 \le  \frac{2}{\eta_x T} \left( \fL_{\lambda}^*(x_0) - \inf_{x \in \BR^{d_x}} \fL_{\lambda}^*(x) + \Vert y_0 - y_{\lambda}^*(x_0) \Vert^2 + \Vert y_0 - y^*(x_0) \Vert^2 \right).
\end{align*}
% The result follows from
% Lemma \ref{lem:1st}, Proposition \ref{prop:iGD} and Proposition \ref{prop:error}.
\end{proof}

Below, we focus on the stochastic setting. We need to use the following theorem for large batch SGD to bound the error from inner loop.

\begin{thm}
\label{thm:SGD-SC}
Suppose $h(x): \BR^d \rightarrow \BR$ is  $\beta$-gradient Lipschitz and $\alpha$-strongly convex. Consider the following update of stochastic gradient descent (SGD):
\begin{align*}
    x_{t+1} = x_t - \frac{1}{\beta} \nabla h(x_t;\fB_t),
\end{align*}
where the mini-batch gradient satisfies 
\begin{align*}
    \BE_{\fB_t} \left[ \nabla h(x_t;\fB_t) \right] = \nabla h(x_t) , \quad \BE_{\fB_t} \Vert \nabla h(x_t;\fB_t) - \nabla h(x_t) \Vert^2 \le \frac{\sigma^2}{B}.
\end{align*}
Then it holds that
\begin{align*}
    \BE \left[ \Vert x_T - x^* \Vert^2 \right] \le \left(1 - \frac{\alpha}{\beta} \right)^T \frac{\beta}{\alpha} \Vert x_0 - x^* \Vert^2 + \frac{\sigma^2}{\alpha^2 B},
\end{align*}
where $x^* = \arg \min_{x \in \BR^d} h(x)$.
\end{thm} 

\begin{proof}
We have that
    \begin{align*}
         \BE \left[ h(x_{t+1}) - h^*\right] 
        &\le \BE \left[ h(x_t) - h^* + \nabla h(x_t)^\top (x_{t+1} - x_t) + \frac{\beta}{2} \Vert x_{t+1} - x_t \Vert^2 \right]\\
        &= \BE \left[ h(x_t) - h^* - \frac{1}{2\beta} \Vert \nabla h(x_t) \Vert^2 + \frac{1}{2 \beta} \Vert \nabla h(x_t) -  \nabla h(x_t; \fB_t) \Vert^2 \right] \\
        &\le \left( 1 - \frac{\alpha}{\beta}\right) \left( h(x_t) - h^* \right) + \frac{\sigma^2}{2 \beta B}.
    \end{align*}
Telescoping gives 
\begin{align*}
    \BE \left[ h(x_T) - h^*\right] \le \left(1 - \frac{\alpha}{\beta} \right)^T  \left( h(x_0) - h^* \right)  + \frac{\sigma^2}{2 \alpha B}.
\end{align*}
We conclude the proof by using the fact that
\begin{align*}
    \frac{\alpha}{2} \Vert x - x^* \Vert^2 \le h(x) - h^* \le \frac{\beta}{2} \Vert x - x^* \Vert^2. 
\end{align*}
\end{proof}
% \begin{proof}
% We have that
%     \begin{align*}
%         &\quad \BE \left[ h(x_{t+1}) - h^*\right] \\
%         &\le \BE \left[ h(x_t) - h^* + \nabla h(x_t)^\top (x_{t+1} - x_t) + \frac{\beta}{2} \Vert x_{t+1} - x_t \Vert^2 \right]\\
%         &= \BE \left[ h(x_t) - h^* - \eta \nabla h(x_t)^\top \nabla h(x_t; \fB_t) + \frac{\beta \eta^2}{2} \Vert \nabla h(x_t; \fB_t) \Vert^2 \right] \\
%         &\le \BE \left[ h(x_t) - h^* - ( \eta - \beta \eta^2) \Vert \nabla h(x_t) \Vert^2  + \beta \eta^2 \Vert \nabla h(x_t; \fB_t) \Vert^2 \right] \\
%         &= \BE \left[ h(x_t) - h^* - \frac{\eta}{2} \Vert \nabla h(x_t) \Vert^2 + \beta \eta^2 \Vert \nabla h(x_t) -  \nabla h(x_t; \fB_t) \Vert^2 \right] \\
%         &\le \left( 1 - \alpha \eta \right) \left( h(x_t) - h^* \right) + \frac{\beta \eta^2 \sigma^2}{B}.
%     \end{align*}
% Telescoping gives 
% \begin{align*}
%     \BE \left[ h(x_T) - h^*\right] \le \left(1 - \alpha \eta \right)^T  \left( h(x_0) - h^* \right)  + \frac{\beta \eta \sigma^2}{\alpha B}.
% \end{align*}
% We conclude the proof by using the fact that
% \begin{align*}
%     \frac{\alpha}{2} \Vert x - x^* \Vert^2 \le h(x) - h^* \le \frac{\beta}{2} \Vert x - x^* \Vert^2. 
% \end{align*}
% \end{proof}

\thmFBSA*

\begin{proof} By taking expectation in \Eqref{eq:des}, we get
\begin{align} \label{eq:des-stoc}
\begin{split}
    &\quad \BE[ \fL_{\lambda}^*(x_{t+1})]  \\
    &\le \BE \left[\fL_{\lambda}^*(x_t) 
    - \frac{\eta_x}{2} \Vert \nabla \fL_{\lambda}^*(x_t) \Vert^2 - 
    \frac{1}{4\eta_x} \Vert x_{t+1} - x_{t} \Vert^2
    + \frac{\eta_x}{2} \Vert G_t - \nabla \fL_{\lambda}^*(x_t) \Vert^2 \right].
\end{split}
\end{align}

Then to prove that this algorithm can output a point such that $\BE \Vert \nabla \varphi(x) \Vert \le \epsilon$, it
suffices to show $\BE \Vert G_t - \nabla \fL_{\lambda}^*(x_t) \Vert^2 = \fO(\epsilon^2)$ holds for all $t$. This requires
\begin{enumerate}[label=\alph*.]
\item $\BE \Vert \nabla_x f(x_t,y_t^K;B_{\rm out}) - \nabla_x f(x_t,y_{\lambda}^*(x_t)) \Vert^2 = \fO(\epsilon^2)$.
\item $\BE \Vert \nabla_x g(x_t,y_t^K;B_{\rm out}) - \nabla_x g(x_t,y_{\lambda}^*(x_t)) \Vert^2 = \fO\left( \epsilon^2/ \lambda^2 \right)$.
\item $  \BE \Vert \nabla_x g(x_t,z_t^K;B_{\rm out}) - \nabla_x g(x_t,y^*(x_t)) \Vert^2 = \fO\left( {\epsilon^2}/{\lambda^2} \right)$.
\end{enumerate}
Our setting of $B_{\rm out}$ shows that it suffices to have
\begin{align*}
\max \left\{ \BE \Vert y_t^K - y_{\lambda}^*(x_t) \Vert^2 ,\BE \Vert z_t^K - y^*(x_t) \Vert^2 \right\}  &= \fO \left( \frac{\epsilon^2}{L_f^2+ \lambda^2 L_g^2} \right), 
% &= \fO \left( \frac{\epsilon^2}{L_f^2 + \lambda^2 L_g^2} \right).
\end{align*}
And our setting of $B_{\rm in}$ and $K_t$ fulfills these conditions if $\Vert y_t^0 - y_{\lambda}^*(x_t) \Vert^2 + \Vert z_t^0 - y^*(x_t) \Vert^2 \le \delta_t$.
Now we use Theorem \ref{thm:SGD-SC} to establish the recursion of $\delta_t$. Note that
\begin{align} \label{eq:recursion-delta-t}
\begin{split}
    &\quad \BE \Vert y_{t+1}^0- y_{\lambda}^*(x_{t+1}) \Vert^2  \\
    & \le 2 \BE \Vert y_t^K- y_{\lambda}^*(x_{t}) \Vert^2 + 2 \BE \Vert y_{\lambda}^*(x_{t}) - y_{\lambda}^*(x_{t+1}) \Vert^2 \\
    &\le \frac{4 L_g}{\mu} \exp \left( - \frac{\mu K_t}{2L_g} \right) \BE \Vert y_t^0 - y_{\lambda}^*(x_t) \Vert^2  + \frac{32 L_g^2}{\mu^2}  \Vert x_{t+1} - x_t \Vert^2 + \frac{ \sigma_g^2}{4 L_g B_{\rm in}} \\
    &\le \frac{1}{2} \BE \Vert y_t^0 - y_{\lambda}^*(x_t) \Vert^2  + \frac{32 L_g^2}{\mu^2}  \Vert x_{t+1} - x_t \Vert^2 + \frac{ \sigma_g^2}{4 L_g B_{\rm in}} 
\end{split}
\end{align}
where the second last inequality uses Theorem \ref{thm:SGD-SC} and the last inequality holds once we get $K_t = \Omega( \kappa \log (\kappa))$.  A similar bound also holds for $\BE \Vert z_{t+1}^0 - y^*(x_{t+1}) \Vert^2$. Putting them together, we get that
\begin{align*}
    \delta_{t+1} \le \frac{1}{2} \delta_t +  \frac{34 L_g^2}{\mu^2}  \Vert x_{t+1} - x_t \Vert^2 + \frac{ \sigma_g^2}{2 L_g B_{\rm in}}
\end{align*}
Then we telescope the recursion of $\delta_t$ and take expectation to obtain that
\begin{align} \label{eq:ub-deltat}
\begin{split}
        \sum_{t=0}^{T-1} \BE[\delta_t] &\le 2 \delta_0 + \frac{68  L_g^2}{\mu^2} \sum_{t=0}^{T-1} \BE [\Vert x_{t+1} - x_t \Vert^2] + \frac{ \sigma_g^2}{ L_g B_{\rm in}} \\
    &=  2 \delta_0 + \fO(\eta_x \Delta + T \eta_x^2 \epsilon^2) + \frac{ \sigma_g^2}{ L_g B_{\rm in}},
\end{split}
\end{align}
where we use \Eqref{eq:des-stoc} in the last step. Let $C$ be a constant that  contains logarithmic factors. We can the bound the expected total iteration number via
\begin{align*}
    \sum_{t=0}^{T-1} \BE[K_t] = \frac{C L_g}{\mu} \sum_{t=0}^{T-1} \BE [\log (\delta_t)]  \le \frac{C L_g T }{\mu} \log \left( \frac{\sum_{t=0}^{T-1} \BE[\delta_t]}{T} \right),
\end{align*}
where we use Jensen's inequality and the concavity of $\log (\,\cdot\,)$. Note that $\sum_{t=0}^{T-1} \BE[\delta_t]$ can be upper bounded via \Eqref{eq:ub-deltat}.
Finally, the total sample complexity of the algorithm is
\begin{align*}
    B_{\rm in} \sum_{t=0}^{T-1} K_t + T B_{\rm out},
\end{align*}
which is the claimed complexity by plugging the choice of $T,B_{\rm in}, B_{\rm out}$.
\end{proof}

% The above theorem requires large batch size $B_{\rm out} \asymp (\sigma_f^2 + \lambda^2 \sigma_g^2) / \epsilon^2$ to make the stochastic algorithm mimic the behavior of the deterministic algorithm. Below, we prove an improved result that allows $\fO(1)$ batch sizes. We prove this result by applying the recently proposed fast two-time-scale stochastic gradient method \citep{zeng2024fast} to solve the Lagrangian function:
% \begin{align*}
%     \min_{x \in \BR^{d_x}} \max_{y \in \BR^{d_y},z \in \BR^{d_y}} \left\{ \fL_{\lambda}(x,y, z) : = f(x,y) + \lambda (g(x,y) - g(x,z) \right\}
% \end{align*}
% We can 

% \begin{align*}
%     \min_x h(x), \quad {\rm where} \quad h(x) = f(x,y^*(x)). 
% \end{align*}

{
\section{Removing the Large Batches in F${}^2$BSA} \label{apx:single-batch}

Theorem \ref{thm:F2BSA} requires large batch size $B_{\rm out} \asymp (\sigma_f^2 + \lambda^2 \sigma_g^2) / \epsilon^2$ to track the deterministic algorithm. Below, we 
get rid of the large batches via a more intricate algorithm design. 
We consider the \textit{single batch scenario}, where
each SGD update samples a single random index $(\phi_f, \phi_g)$. Define the following stochastic gradient
\begin{align}
\begin{split}
    &\quad \nabla \fL_{\lambda}^*(x; \phi_f, \phi_g) =  \nabla \fL_{\lambda}(x, y_{\lambda}^*(x); \phi_f, \phi_g), \\
    &{\rm where}~~ \nabla \fL_{\lambda}(x, y; \phi_f, \phi_g) = \nabla_x f(x,y; \phi_f) + \lambda (\nabla_x g(x,y; \phi_g) - \nabla g^*(x; \phi_g)).
\end{split}
\end{align}
It is clear that $\BE \nabla \fL_{\lambda}^*(x; \phi_f, \phi_g) = \nabla \fL_{\lambda}^*(x)$. 
Our improved convergence rate in Theorem \ref{thm:F2BSA} requires the following condition on the hyper-gradient estimator $G$.

\begin{algorithm*}[t]  
\caption{F${}^3$BSA$(x, y_0)$} \label{alg:F3BSA}
\begin{algorithmic}[1]
\STATE $z_0 = y_0$ \\[1mm]
\STATE \textbf{for} $ t=0,1,\cdots,T-1$ \\[1mm]
\STATE \quad Sample random index $(\phi_f^t, \phi_g^t)$ \\[1mm]
\STATE \quad $(G_t, y_{t+1} , z_{t+1}) = \text{MLMC-HyperGradEst}(x_t, y_t, z_t, \phi_f^t, \phi_g^t, M_t, N_t, S_t )$ \\[1mm]
\STATE \quad $x_{t+1} = x_t - \eta_x G_t$ \\[1mm]
\STATE \textbf{end for} 
\end{algorithmic}
\end{algorithm*}
\begin{algorithm*}[t]  

\caption{MLMC-HyperGradEst$(x, y_0, z_0, \phi_f, \phi_g, M,  N,  S)$} \label{alg:F2HE}
\begin{algorithmic}[1] 
\STATE \textbf{for} $m=1,\cdots,M$ \\ [1mm]
\STATE \quad Draw $J \sim {\rm Geom}(1/2)$ \\[1mm]  
\STATE \quad Write $(G_m^{(j)}, y_m^{(j)} , z_m^{(j)}) = \text{HyperGradEst}(x,y_0,z_0, N, j )$ \\[1mm]
\STATE \quad $G_m = G_m^{(0)} + 2^J (G_m^{(J)} - G_m^{(J-1)}) \BI[ J \le S] $ \\[1mm] % \STATE \quad $g_m^N = $ \\[1mm]
% \STATE \quad $y_m = y_m^{(0)} + 2^J (y_m^{(J)} - y_m^{(J-1)}) \BI[ J \le J_{\max}] $ \\[1mm]
% \STATE \quad $z_m = z_m^{(0)} + 2^J (z_m^{(J)} - z_m^{(J-1)}) \BI[ J \le J_{\max}] $ \\[1mm]
% \STATE \quad \textbf{if} $J \le J_{\max} $ \\[1mm] 
% \STATE \quad \quad $g_m = g_m^N + 2^J \left( \text{SingleF${}^2$HE}(x,y_0,z_0, N, J ) - 
% \text{SingleF${}^2$HE}(x,y_0,z_0, N, J-1 )
% \right) $ \\[1mm]
\STATE \textbf{end for} \\[1mm]
\STATE $ \hat G = \frac{1}{M} \sum_{m=1}^M G_m$, $ \hat y = y_1^{(0)}$, $ \hat z = y_1^{(0)}$ \\[1mm]
\STATE \textbf{return} $(\hat G, \hat y , \hat z)$ 
\end{algorithmic}
\end{algorithm*}

\begin{algorithm*}[t]  
\caption{HyperGradEst$(x, y_0, z_0, \phi_f, \phi_g,  N, K)$} \label{alg:single-F2HE}
\begin{algorithmic}[1] 
\STATE $\hat y = \text{SGD}{}^{\rm SC}( f(x,\,\cdot\,; \phi_f) + \lambda g(x, \,\cdot\,; \phi_g), y_0, N, K) $ \\[1mm]
\STATE $\hat z = \text{SGD}{}^{\rm SC}( \lambda g(x,\,\cdot\,; \phi_g) , z_0, N, K) $ \\[1mm]
\STATE $\hat G = \nabla_x f(x,\hat y; \phi_f) + \lambda (\nabla_x g(x, \hat y; \phi_g) - \nabla_x g(x, \hat z; \phi_g)) $ \\[1mm]
\STATE \textbf{return} $(\hat G, \hat y, \hat z)$
\end{algorithmic}
\end{algorithm*}

\begin{asm} \label{asm:hyper-bias-variance}
Assume that the hyper-gradient estimator $G$ satisfies 
\begin{align} \label{eq:cond-HE}
\begin{split}
 \Vert \BE G -  \nabla \fL_{\lambda}^*(x; \phi_f, \phi_g )  \Vert^2 \le \zeta_{\rm bias} &= \fO( \min\{ \epsilon^2, \sigma_f^2 + \lambda^2 \sigma_g^2 \}), \\
\BE \Vert G - \BE G \Vert^2 \le \zeta_{\rm varinace} &= \fO(\sigma_f^2 + \lambda^2 \sigma_g^2),
\end{split}
\end{align}
\end{asm}
\begin{restatable}{thm}{thmouter} \label{thm:outer}
    Suppose Assumption \ref{asm:sc} and \ref{asm:stochastic} hold. 
Let $\ell$,  $\kappa$ defined as Definition \ref{dfn:kappa-sc}. Let
$\lambda \asymp \max\{ \ell \kappa^2 / \Delta, ~\ell \kappa^3 / \epsilon \}$, where $\Delta:=\varphi(x_0) - \inf_{x \in \BR^{d_x}} \varphi(x)$.
If for any $t$ the hyper-gradient estimator $G_t$ satisfies Assumption \ref{asm:hyper-bias-variance},
then iterating $ x_{t+1}  = x_t - \eta_x G_t$ with
\begin{align} \label{eq:eta-outer-F2BSA}
    \eta_x \asymp \frac{\epsilon^2}{(\sigma_f^2 + \lambda^2 \sigma_g^2)  \ell 
 \kappa^3}, \quad T \asymp \frac{\Delta}{\eta_x \epsilon^2} 
\end{align}
output $\bar x_T = \frac{1}{T} \sum_{t=0}^{T-1} x_t$ such that $\BE \Vert \nabla \varphi(\bar x_T) \Vert \le \epsilon$. 
\end{restatable}




In Assumption \ref{asm:hyper-bias-variance}, a nearly unbiased hyper-gradient estimator $G_t$ is required to recover the same convergence rate as SGD on $\varphi(x)$. In F${}^2$BSA (Algorithm \ref{alg:F2BSA}), this condition of ${\rm bias} \lesssim \epsilon^2$ is achieved by running inner-loop SGD with large batch to accuracy $ \epsilon^2$. In the following, we show it is possible to use existing technique~\citep{asi2021stochastic,hu2021bias} for improving biased SGD, which helps relax the accuracy of inner-loop SGD to $ \sigma_f^2 + \lambda^2 \sigma_g^2$. In this case, using single-batch SGD would not harm the convergence rate. We call this improved algorithm F${}^3$BSA (see Algorithm \ref{alg:F3BSA}). The ultimate algorithms look a little complicated at first glance, but they all come from standard techniques in the literature. We summarize the key steps below.
\begin{enumerate}
    \item F${}^3$BSA (Algorithm \ref{alg:F3BSA}) conducts a SGD update on variable $x$ with the MLMC hyper-gradient estimator $G_t$ given by Algorithm \ref{alg:F2HE}.
    \item The design of MCMC-HyperGradEst (Algorithm \ref{alg:F2HE}) has the same structure as the standard MLMC estimator \citep{giles2008multilevel,blanchet2015unbiased}. Following \citep{asi2021stochastic,hu2021bias}, we wrap a high-bias hyper-gradient estimator (Algorithm~\ref{alg:single-F2HE}) into the MLMC structure to return a low-bias estimator. 
    \item HyperGradEst (Algorithm \ref{alg:single-F2HE}) applies SGD updates to solve the lower-level problems with respect to variables $y$ and $z$ in \Eqref{eq:nabla-Lag}. 
    Since the naive SGD update in F${}^2$BSA can not meet the 
    requirements in MLMC, we use the SGD${}^{
    \rm SC}$ algorithm \citep[Algorithm 2]{allen2018make}, which is presented in Section \ref{apx:SGD-SC} for completeness. It runs multiple epochs of SGD with different step sizes like \citep{hazan2014beyond} to achieve the optimal rate in smooth strongly-convex stochastic first-order optimization.
\end{enumerate}


The following theorem shows Algorithm \ref{alg:F2HE} can output a nearly unbiased hyper-gradient estimator at a low cost. The proof follows from \citep{asi2021stochastic}.

\begin{restatable}{thm}{thminner} \label{thm:inner}
Suppose Assumption \ref{asm:sc} and \ref{asm:stochastic} hold. 
Let $\ell$,  $\kappa$ defined as Definition \ref{dfn:kappa-sc}. For any given $x \in \BR^{d_x}$, then Algorithm \ref{alg:F2HE} with 
\begin{align} \label{eq:hyper-F2HE}
    N_t = S_t \asymp \log \left( \frac{\max\{ \lambda^2 \ell^2 \delta_t,   \kappa (\sigma_f^2 + \lambda^2 \sigma_g^2) \}}{ \zeta_{\rm bias} } \right), \quad M_t \asymp \frac{\kappa (\sigma_f^2 + \lambda^2 \sigma_g^2) N_t}{\zeta_{\rm variance}}
\end{align}
outputs an hyper-gradient estimator $G_t$ satisfying Condition \ref{eq:cond-HE} in $\fO( \kappa N_t M_t)$ SFO complexity, where $\delta_t$ is the upper bound of $\Vert y_t^0 - y_{\lambda}^*(x_t) \Vert^2 + \Vert z_t^0 - y^*(x_t) \Vert^2$.
% \begin{align*}
%     \Vert \BE \hat G - \nabla \fL_{\lambda}^*(x) \Vert^2 \le \zeta_{\rm bias}, \quad \BE \Vert \hat G - \BE \hat G \Vert^2 \le \zeta_{\rm variance},
% \end{align*}
\end{restatable}

Combining Theorem \ref{thm:outer} and \ref{thm:inner}, we obtain the complexity of Algorithm \ref{alg:F3BSA} as follows. 

\begin{restatable}{thm}{thmFSBAsingle} \label{thm:F3BSA}
Suppose Assumption \ref{asm:sc} and \ref{asm:stochastic} hold. 
Let $\ell$,  $\kappa$ defined as Definition \ref{dfn:kappa-sc}.
Define
$\Delta:=\varphi(x_0) - \inf_{x \in \BR^{d_x}} \varphi(x)$ and $R := \Vert y_0 - y^*(x_0) \Vert^2 $.
Let $\lambda \asymp \max\left\{\kappa/R,~ \ell \kappa^2/ \Delta, ~ \ell \kappa^3 / \epsilon \right\} $ 
and set other parameters in Algorithm \ref{alg:F3BSA} according to \Eqref{eq:eta-outer-F2BSA} and \Eqref{eq:hyper-F2HE}, 
where $\delta_t$ is defined via the recursion
\begin{align} \label{eq:recursion-delta-single}
    \delta_{t+1} = \frac{1}{2}\delta_t + \frac{34 L_g^2}{\mu^2}  \Vert x_{t+1} - x_t \Vert^2 + \frac{24 (\sigma_f^2+ \lambda^2 \sigma_g^2)}{\lambda \mu (L_f + \lambda L_g)}, \quad \delta_0 = \fO(R),
\end{align}
then it
can output a point such that $\BE \Vert \nabla \varphi(x) \Vert \le \epsilon$ in the total SFO complexity of
\begin{align*}
\begin{cases}
    \fO(\ell \kappa^5 \epsilon^{-4} \log (\nicefrac{ \ell \kappa}{\epsilon})), & \sigma_f >0 , \sigma_g = 0;\\
    {\fO(\ell^3 \kappa^{11} \epsilon^{-6} \log (\nicefrac{ \ell \kappa}{\epsilon})) }, & \sigma_f>0, \sigma_g >0.
\end{cases}
\end{align*}
\end{restatable}

The above theorem shows that F${}^3$BSA achieves the $\tilde \fO(\epsilon^{-4})$ and $\tilde \fO(\epsilon^{-6})$ SFO complexity for the partially and fully stochastic cases as F${}^2$BSA while using a single batch at each iteration. 
Another side advantage of F${}^3$BSA is that it also improves the dependency in $\kappa$ because its inner loop uses the optimal stochastic algorithm SGD${}^{\rm SC}$ \citep[Algorithm 2]{allen2018make} to achieve better bias-variance tradeoff, instead of using naive SGD in F${}^2$BSA.


% Let $\eta_x \le 1/(2L)$ and telescoping yields
% \begin{align} \label{eq:SGD-bias-variance}
%     \frac{1}{T} \sum_{t=0}^{T-1} \BE \Vert \nabla \fL_{\lambda}^*(x_t) \Vert^2 \le \frac{4 (\fL_{\lambda}^*(x_t) - \inf_{x \in \BR^{d_x}} \fL_{\lambda}^*(x))}{ \eta_x T} + \underbrace{2 \times \zeta_{\rm bias} + 2 \eta_x L \times \zeta_{\rm variance}}_{\rm error}.
% \end{align}
% To guarantee an $\epsilon$-stationary point, we need the error term be $\fO(\epsilon^2)$. 
% It reveals the challeng for using small batch in our setting. For standard SGD analysis on $\fL_{\lambda}^*(x)$, it is assumed that $\hat \nabla \fL_{\lambda}^*(x_t)$ is an unbiased and the error only appears in the variance term, which can be controlled by scaling down the step size. In contrast, in our setting, our estimator of $\nabla \fL_{\lambda}^*(x_t)$ is biased due to the inexactness in the inner loop solving, and the error caused by the bias term can not be improved even if we use a smaller step size.

% To tackle the above challenge, we use the bias reduction  \citep{asi2021stochastic} to return a nearly unbiased hyper-gradient estimator.
% The original work of \citep{asi2021stochastic} achieves bias reduction by applying Multilevel Monte-Carlo (MLMC) technique on the EpochSGD algorithm \citep{hazan2014beyond} for Lipschitz objectives. However, we do not assume that $\fL_{\lambda}^*(x)$ is Lipschitz in our setting, and we need to adapt the EpochSGD for smooth objectives as in \citep{allen2018make}. The convergence in each epoch relies on the following standard result for SGD.





% \begin{thm} [{\citet[Theorem 4.1 (b)]{allen2018make}}]
% Suppose $h(x): \BR^d \rightarrow \BR$ is  $\beta$-gradient Lipschitz and $\alpha$-strongly convex. Epoch SGD algorithm outputs $x_{\rm out}$ satisfying
% \begin{align*}
% % \BE \Vert x_N - x^* \Vert^2 &\le \frac{2 \sigma^2}{\alpha \beta} + \frac{2 \Vert x_0 - x^* \Vert^2}{2^{ K \alpha / (8 \beta)}}   \\
% \BE \Vert x_{\rm out} - x^* \Vert^2 &\le  \frac{2\sigma^2}{\alpha^2 K} + \frac{2 \Vert x_0 - x^* \Vert^2}{2^{ N + 2K}}?
%     % \BE \Vert x_{\rm out} - x^* \Vert^2 \le  \fO \left(\frac{\sigma^2}{\alpha^2 K} \right)  + 2^{- \alpha T / (8 \beta)} \Vert x_0 - x^* \Vert^2.
% \end{align*}
% \end{thm}
% Apply Epoch SGD for smooth objectives in variables $y$ and $z$. We obtain that 
% \begin{align*}
%     \BE \Vert g_t^K - \nabla \fL_{\lambda}^*(x_t) \Vert^2 \le \frac{1}{L_f^2 + \lambda^2 L_g^2}  \cdot \left( \frac{\sigma^2 C}{\alpha^2 K}  + \left(1 - \frac{\alpha}{\beta} \right)^{K / C} \Vert y_t - y_{\lambda}^*(x_t) \Vert^2 + \Vert z_t - y^*(x_t) \Vert^2 \right).
% \end{align*}

% \begin{thm}
% Suppose Assumption \ref{asm:sc} and \ref{asm:stochastic} hold. 
% Let $\ell$,  $\kappa$ defined as Definition \ref{dfn:kappa-sc}. For any given $x \in \BR^{d_x}$, Algorithm \ref{alg:F2HE} with 
% \begin{align} \label{eq:hyper-F2HE}
%     N = J_{\max} \asymp \log \left( \frac{\max\{ \lambda^2 \ell^2 \delta,   \kappa (\sigma_f^2 + \lambda^2 \sigma_g^2) \}}{ \zeta_{\rm bias} } \right), \quad M \asymp \frac{\kappa (\sigma_f^2 + \lambda^2 \sigma_g^2) J_{\max}}{\zeta_{\rm variance}}.
% \end{align}
% outputs an hyper-gradient estimator $\hat g$ such that
% \begin{align*}
%     \Vert \BE \hat g - \nabla \fL_{\lambda}^*(x) \Vert^2 \le \zeta_{\rm bias}, \quad \BE \Vert \hat g - \BE \hat g \Vert^2 \le \zeta_{\rm variance},
% \end{align*}
% where $\Vert y_0 - y_{\lambda}^*(x) \Vert^2 + \Vert z_0 - y^*(x) \Vert^2 \le \delta$. The total SFO complexity is $\tilde \fO( \kappa M)$.
% \end{thm}

\subsection{Missing Proofs for F${}^3$BSA} \label{apx:proof-single-batch}

We present the missing proofs in this section. 

\thmouter* 
\begin{proof} Let $L$ be the constant of gradient Lipschitz continuity of $\fL_{\lambda}^*(x)$. From 
Lemma \ref{lem:1st} we know that $L = \fO(\ell \kappa^3)$. We have that
    \begin{align*}
    \BE[ \fL_{\lambda}^*(x_{t+1})] &\le \BE \left[\fL_{\lambda}^*(x_t) + \langle \nabla \fL_{\lambda}^*(x_t) , x_{t+1} - x_t
    \rangle + \frac{L}{2} \Vert x_{t+1 } - x_t \Vert^2 \right] \\
    &= \BE \left[ \fL_{\lambda}^*(x_t) -
    \eta_x \nabla \fL_{\lambda}^*(x_t)^\top G_t + \frac{\eta_x^2 L}{2} \Vert G_t \Vert^2 \right] \\
    &\le \fL_{\lambda}^*(x_t) - \left( \frac{\eta_x}{2} - \frac{\eta_x^2 L}{2} \right) \Vert \nabla \fL_{\lambda}^*(x_t) \Vert^2  \\
    &\quad + \frac{\eta_x}{2} \Vert \BE G_t -  \nabla \fL_{\lambda}^*(x_t; \phi_f^t, \phi_g^t )  \Vert^2 + \frac{\eta_x^2 L}{2} \BE \Vert G_t -  \nabla \fL_{\lambda}^*(x_t)  \Vert^2 \\
    &\le \fL_{\lambda}^*(x_t) - \left( \frac{\eta_x}{2} - \frac{\eta_x^2 L}{2} \right) \Vert \nabla \fL_{\lambda}^*(x_t) \Vert^2 + \frac{\eta_x}{2} \Vert \BE G_t -  \nabla \fL_{\lambda}^*(x_t; \phi_f^t, \phi_g^t )  \Vert^2 \\
    &\quad  + \eta_x^2 L ( \BE \Vert G_t -  \nabla \fL_{\lambda}^*(x_t; \phi_f^t, \phi_g^t)  \Vert^2 + \BE \Vert \nabla \fL_{\lambda}^*(x_t; \phi_f^t, \phi_g^t) - \nabla \fL_{\lambda}^*(x_t; \phi_f^t, \phi_g^t) \Vert^2 ).
    % &\le \fL_{\lambda}^*(x_t) - \left( \frac{\eta_x}{2} - \frac{\eta_x^2 L}{2} \right) \Vert \nabla \fL_{\lambda}^*(x_t) \Vert^2  \\
    % &\quad + \frac{\eta_x}{2} \underbrace{\Vert \BE g_t -  \nabla \fL_{\lambda}^*(x_t)  \Vert^2}_{\le {\zeta_{\rm bias}}} + \frac{\eta_x^2 L}{2} \underbrace{\BE \Vert g_t -  \nabla \fL_{\lambda}^*(x_t)  \Vert^2}_{\le \zeta_{\rm bias} +  \zeta_{\rm variance}}.
    % \left( \frac{\eta_x}{2} - \frac{\eta_x^2 L}{2} \right) \Vert \hat \nabla \fL_{\lambda}^*(x_t) \Vert^2 + \frac{\eta_x}{2} \Vert \hat  \nabla \fL_{\lambda}^*(x_t) - \nabla \fL_{\lambda}^*(x_t) \Vert^2 \right] \\
    \end{align*}
    Plugging in \Eqref{eq:cond-HE} and \Eqref{eq:eta-outer-F2BSA}, we have that
\begin{align*}
     \BE[ \fL_{\lambda}^*(x_{t+1})] \le \fL_{\lambda}^*(x_t) - \frac{\eta_x}{4} \Vert \nabla \fL_{\lambda}^*(x_t) \Vert^2  + \fO(\eta_x \epsilon^2).
\end{align*}
Telescope over $t =0,1,\cdots,T-1$, we get
\begin{align} \label{eq:final-fosp}
    \frac{1}{T} \sum_{t=0}^{T-1} \BE \Vert \nabla \fL_{\lambda}^*(x_t) \Vert^2 \le \frac{4 (\fL_{\lambda}^*(x_t) - \inf_{x \in \BR^{d_x}} \fL_{\lambda}^*(x))}{ \eta_x T} + \fO(\epsilon^2),
\end{align}
which outputs an $\epsilon$-stationary point in $T = \fO( \Delta / (\eta_x \epsilon^2) )$ iterations.
\end{proof} 

\thminner* 
\begin{proof} % Define
% \begin{align*}
%     &\quad \nabla \fL_{\lambda}^*(x; \phi_f, \phi_g) =  \nabla \fL_{\lambda}(x, y_{\lambda}^*(x); \phi_f, \phi_g), \\
%     {\rm where} &~~ \nabla \fL_{\lambda}(x, y; \phi_f, \phi_g) = \nabla_x f(x,y; \phi_f) + \lambda (\nabla_x g(x,y; \phi_g) - \nabla g^*(x; \phi_g)).
% \end{align*}
By applying Theorem \ref{thm:SGD-SC}, we know that Algorithm \ref{alg:single-F2HE} ensures that
\begin{align} \label{eq:single-F2HE}
\begin{split}
    \BE \Vert G_m^{(j)} - \nabla \fL_{\lambda}^*(x_t) \Vert^2 & \le (L_f^2 + \lambda^2 L_g^2) \cdot \left( \Vert y_m^{(j)} - y_{\lambda}^*(x_t) \Vert^2 + \Vert z_m^{(j)} - y^*(x_t) \Vert^2 \right)  \\
    &\le (L_f^2 + \lambda^2 L_g^2)
    \cdot \left( \frac{\delta_t}{2^{N+2K}} + \frac{\sigma_f^2 + \lambda^2 \sigma_g^2}{2^{K-3} \lambda \mu (L_f + \lambda L_g)} \right),
\end{split}
\end{align}
The fact $\BP(J=j)=2^{-j}$ implies that 
\begin{align*} 
    \BE \hat G = \BE G_1 = \BE G_1^{(0)} + \sum_{j=1}^{S_t} \BP(J=j) 2^j ( G_1^{(j)} - G_1^{(j-1)}  ) = \BE \left[ G_1^{S_t} \right].
\end{align*}
Therefore, using \Eqref{eq:single-F2HE} the bias of our estimator can be upper bounded by 
\begin{align*}
    &\quad \Vert \BE \hat G - \nabla \fL_{\lambda}^*(x_t) \Vert^2 = \Vert \BE G_1^{S_t} - \nabla \fL_{\lambda}^*(x_t) \Vert^2 \le \BE \Vert G_1^{S_t} - \nabla \fL_{\lambda}^*(x_t) \Vert^2 \\
    &\le (L_f^2 + \lambda^2 L_g^2)\cdot \left( \frac{\delta_t}{2^{N_t +2 S_t}} + \frac{\sigma_f^2 + \lambda^2 \sigma_g^2}{2^{S_t-3} \lambda \mu (L_f + \lambda L_g)  } \right).
\end{align*}
It is easy to check that $\Vert \BE \hat G - \nabla \fL_{\lambda}^*(x) \Vert^2 \le \zeta_{\rm bias}$ holds by our hyper-parameter setting. 
Below, we analyze the variance of our estimator. Note that
\begin{align*}
    &\quad \BE \Vert G_m - G_m^{(0)} \Vert^2 \\
    &= \sum_{j=1}^{S_t} \BP( J = j) 2^{2j} \BE \Vert G_m^{(j)} - G_m^{(j-1)} \Vert^2 \\
    &= \sum_{j=1}^{S_t} 2^j \BE \Vert G_m^{(j)} - G_m^{(j-1)} \Vert^2 \\
    & \le \sum_{j=1}^{S_t} 2^j \cdot \left( 2 \BE \Vert G_m^{(j)} - \nabla \fL_{\lambda}^*(x_t) \Vert^2 + 2 \BE \Vert G_m^{(j-1)} - \nabla \fL_{\lambda}^*(x_t) \Vert^2 \right) \\
    &\le (L_f^2 + \lambda^2 L_g^2) \sum_{j=1}^{S_t} \frac{10 \delta_t}{2^{N+j}} + \frac{48 (\sigma_f^2 + \lambda^2 \sigma_g^2)}{ \lambda \mu (L_f + \lambda L_g)} \\
    &\le (L_f^2 + \lambda^2 L_g^2)  \cdot \left(\frac{10 \delta_t}{2^N} +  \frac{48 (\sigma_f^2 + \lambda^2 \sigma_g^2)}{ \lambda \mu (L_f + \lambda L_g)} \cdot S_t \right),
\end{align*}
where the second last line uses \Eqref{eq:single-F2HE}.
Therefore, we have that
\begin{align*}
    &\quad \BE \Vert \hat G - \BE \hat G \Vert^2 = \frac{1}{M_t} \BE \Vert G_1 - \BE G_1 \Vert^2 \le \frac{1}{M_t} \BE \Vert G_1 - \nabla \fL_{\lambda}^*(x_t) \Vert^2 \\
    &\le \frac{2}{M_t} \BE \Vert G_1 - G_1^{(0)} \Vert^2 + \frac{2}{M_t} \BE \Vert G_1^{(0)} - \nabla \fL_{\lambda}^*(x_t) \Vert^2  \\
    &\le \frac{L_f^2 + \lambda^2 L_g^2}{M_t}  \cdot \left(\frac{22 \delta_t}{2^{N_t}} +  \frac{100 (\sigma_f^2 + \lambda^2 \sigma_g^2)}{ \lambda \mu (L_f + \lambda L_g)} \cdot S_t \right).
\end{align*}
It is also easy to check that $\BE \Vert \hat G - \BE \hat G \Vert^2 \le \zeta_{\rm variance}$ holds by our hyper-parameter setting.
Finally, the expected number of SFO calls is 
\begin{align*}
    \frac{8(L_f + \lambda L_g)}{\lambda \mu} \left( N_t + \sum_{j=1}^{S_t} \BP(J=j)  2^j  \right) M_t  = \frac{8(L_f + \lambda L_g)}{\lambda \mu} \left( N_t + S_t \right) M_t = \fO(\kappa N_t M_t).
\end{align*}
\end{proof}

\thmFSBAsingle* 

\begin{proof}
By Theorem \ref{thm:outer} the outer-loop complexity is $ T =  \fO \left( (\sigma_f^2 + \lambda^2 \sigma_g^2) \ell \kappa^3 \Delta \epsilon^{-4} \right) $. By Theorem \ref{thm:inner} the inner-loop complexity in the $t$-th iteration is $\tilde \fO \left( \kappa^2 \log^2 \delta_t \right)$ for $\delta_t$ satisfying $\Vert y_t^0 - y_{\lambda}^*(x_t) \Vert^2 + \Vert z_t^0 - y^*(x_t) \Vert^2 \le \delta_t$.
Similar to Theorem \ref{thm:F2BSA}, we also require to specify the recursion for $\delta_t$ with known problem parameters. Similar to \Eqref{eq:recursion-delta-t},we have that
\begin{align*}
   \BE \Vert y_{t+1}-  y_{\lambda}^*(x_{t+1}) \Vert^2  &\le 2 \BE \Vert y_{t+1} - y_{\lambda}^*(x_{t}) \Vert^2 + 2 \BE \Vert y_{\lambda}^*(x_{t}) - y_{\lambda}^*(x_{t+1}) \Vert^2 \\
    &\le \frac{1}{2} \Vert y_t - y_{\lambda}^*(x_t) \Vert^2 + \frac{16 (\sigma_f^2 + \lambda^2 \sigma_g^2)}{ \lambda \mu (L_f + \lambda L_g)} + \frac{32 L_g^2}{\mu^2} \Vert x_{t+1} - x_t \Vert^2, 
\end{align*}
where the last step follows Theorem \ref{thm:optimal-SGD-SC} with $N_t \ge 8$ and Lemma \ref{lem:nabla-y-lambda-bound} that $y_{\lambda}^*(x)$ is $(4L_g/\mu)$-Lipschitz continuous. Similarly, we also have that
\begin{align*}
\BE \Vert z_{t+1}-  y^*(x_{t+1}) \Vert^2   &\le 2 \BE \Vert z_{t+1} - y^*(x_{t}) \Vert^2 + 2 \BE \Vert y^*(x_{t}) - y^*(x_{t+1}) \Vert^2 \\
    &\le \frac{1}{2} \Vert z_t - y^*(x_t) \Vert^2 + \frac{8 \sigma_g^2}{ \mu \lambda L_g} + \frac{2 L_g^2}{\mu^2} \Vert x_{t+1} - x_t \Vert^2, 
\end{align*}
where the last step follows Theorem \ref{thm:optimal-SGD-SC} with $N_t \ge 8$
and that $y^*(x)$ is $(L_g/\mu)$-Lipschitz continuous. Combining the above two inequalities yields \Eqref{eq:recursion-delta-single}. Telescoping \Eqref{eq:recursion-delta-single} yields that
\begin{align*}
    \sum_{t=0}^{T-1} \BE[\delta_{t}] &\le 2\delta_0  + \frac{48 (\sigma_f^2+ \lambda^2 \sigma_g^2)}{\lambda \mu (L_f + \lambda L_g)} + \frac{68 L_g^2}{\mu^2} \sum_{t=0}^{T-1} \BE  \Vert x_{t+1} - x_t \Vert^2 \\
    &\le 2\delta_0 + \frac{48 (\sigma_f^2+ \lambda^2 \sigma_g^2)}{\lambda \mu (L_f + \lambda L_g)} + \frac{136 L_g^2 \eta_x^2}{\mu^2} \sum_{t=0}^{T-1} \BE [\Vert 
    \nabla \fL_{\lambda}^*(x_t)
    \Vert^2+ \Vert G_t - \nabla \fL_{\lambda}^*(x_t) \Vert^2] \\
    &\le 2\delta_0 + \frac{48 (\sigma_f^2+ \lambda^2 \sigma_g^2)}{\lambda \mu (L_f + \lambda L_g)} + \frac{136 L_g^2 \eta_x^2 T}{\mu^2} \cdot \left( \epsilon^2 + \sigma_f^2 + \lambda^2 \sigma_g^2 \right).
\end{align*}
Therefore, the total SFO complexity (in expectation) is
\begin{align*}
    \tilde \fO \left( \kappa^2 \sum_{t=0}^{T-1} \log^2 \delta_t \right) \le \tilde \fO \left( \kappa^2 T \max_{ 0 \le t \le T-1} \log^2 \delta_t \right) = \tilde \fO(\kappa^2 T),
\end{align*}
where $T$ is given in \Eqref{eq:eta-outer-F2BSA}.
\end{proof}

\subsection{The SGD Subroutine} \label{apx:SGD-SC}


This section recalls the optimal stochastic algorithm for smooth strongly-convex optimization \citep{allen2018make}, which be used as the subroutines in F${}^3$BSA (Algorithm \ref{alg:F3BSA}). 

\begin{asm} \label{asm:sgd}
Assume that the stochastic gradient satisfies 
\begin{align*}
    \BE_{\phi_t} \left[ \nabla h(x_t;\phi_t) \right] = \nabla h(x_t) , \quad \BE_{\phi_t} \Vert \nabla h(x_t;\phi_t) - \nabla h(x_t) \Vert^2 \le \sigma^2.
\end{align*}
\end{asm}

\begin{algorithm*}[htbp]  
\caption{SGD $(h, x_0, \eta, K)$} \label{alg:SGD-C}
\begin{algorithmic}[1] 
\STATE \textbf{for} $k = 0,1,\cdots,K-1$ \\
\STATE \quad Sample a random index $\phi_t$ \\[1mm]
\STATE \quad $x_{k+1} = x_t - \eta \nabla h(x_t ; \phi_t)$ \\[1mm]
\STATE \textbf{end for} \\[1mm]
\STATE \textbf{return} $\bar x_K = \frac{1}{K} \sum_{k=1}^K x_k$  
\end{algorithmic}
\end{algorithm*}

\begin{thm}{\citet[Theorem 4.1 (a)]{allen2018make}}
\label{thm:SGD-C}
Suppose $h(x): \BR^d \rightarrow \BR$ is  $\beta$-gradient Lipschitz and Assumption \ref{asm:sgd} holds for the stochastic gradient. Algorithm \ref{alg:SGD-C} outputs $\bar x_K$ such that
\begin{align*}
    \BE h(\bar x_K)  - h(x^*) \le \frac{\Vert x_0 -x^* \Vert^2}{2\eta K} +  \frac{\eta \sigma^2}{2 (1- \eta \beta)},
\end{align*}
where $x^* = \arg \min_{x \in \BR^{d_x}} h(x)$.
% strongly convex. 

% Consider the following update of stochastic gradient descent (SGD):
% \begin{align*}
%     x_{t+1} = x_t - \eta \nabla h(x_t;\xi_t),
% \end{align*}
% where
% Then it holds that
% \begin{align*}
%     \BE \left[ \Vert x_T - x^* \Vert^2 \right] \le \left(1 - \frac{\alpha}{\beta} \right)^T \frac{\beta}{\alpha} \Vert x_0 - x^* \Vert^2 + \frac{\sigma^2}{\alpha^2 B},
% \end{align*}
% where $x^* = \arg \min_{x \in \BR^d} h(x)$.
\end{thm} 


Based on the above result, we can analyze the SGD${}^{\rm SC}$ algorithm \citep[Algorithm 2]{allen2018make}, which achieves the optimal rate for smooth strongly-convex stochastic optimization. We present this algorithm in Algorithm \ref{alg:SGD-SC}. The convergence of Algorithm \ref{alg:SGD-SC} can be found in {\citep[Theorem 4.1 (b)]{allen2018make}}. However, they did not explicitly calculate the constants in the convergence rates. Therefore, we provide a self-contained proof below.

\begin{algorithm*}[htbp]  
\caption{SGD${}^{\rm SC}$ $(h, x_0, N, K)$} \label{alg:SGD-SC}
\begin{algorithmic}[1] 
\STATE \textbf{for} $k = 1,\cdots,N$ \quad \textbf{do} $x_{k} = \text{SGD}(h, x_{k-1}, 1/(2\beta), 4 \beta /\alpha) $ \\ [1mm]
\STATE \textbf{for} $k= 1,\cdots,K$ \quad \textbf{do} $x_{N+k} = \text{SGD}(h, x_{N+k-1},  1/ (2^k \beta), 2^{k+2} \beta /\alpha)$ \\ [1mm] 
\STATE \textbf{return} $x_{N+K}$   
\end{algorithmic}
\end{algorithm*}

\begin{thm} \label{thm:optimal-SGD-SC}
    Suppose $h(x): \BR^d \rightarrow \BR$ is $\beta$-gradient Lipschitz and $\alpha$-strongly convex. Algorithm \ref{alg:SGD-SC} outputs $x_{N+K}$ satisfying
    \begin{align}
        \BE \Vert x_{N+K} - x^* \Vert^2 &\le \frac{\Vert x_0 - x^* \Vert^2}{2^{N+2K}} + \frac{\sigma^2}{2^{K-2} \alpha \beta}.
        \label{eq:EpochSGD-xK}
    \end{align}
The total stochastic first-order oracle (SFO) complexity is bounded by $ \lceil 4 \kappa ( N + 2^K) \rceil $, where $\kappa = \beta /\alpha$ is the condition number.
\end{thm}

\begin{proof}
Under the strong convexity assumption, Theorem \ref{thm:SGD-C} yields that
\begin{align} \label{eq:each-epoch}
    \BE \Vert \bar x_K - x^* \Vert^2 \le \frac{2}{\alpha} [\BE h(\bar x_K)  - h(x^*)] \le  \frac{\Vert x_0 - x^* \Vert^2}{\alpha \eta K} + \frac{\eta \sigma^2}{\alpha (1 - \eta \beta)}.
\end{align}
For the first $N$ epochs,  by the parameter setting in Algorithm \ref{alg:SGD-SC} and \Eqref{eq:each-epoch}, we have
\begin{align*}
    \BE \Vert x_{k} - x^* \Vert^2 \le \frac{\Vert x_{k-1} - x^* \Vert^2}{2} + \frac{\sigma^2}{\alpha \beta},
\end{align*}
Telescoping for all the $N$ epochs yields
\begin{align}
     \BE \Vert x_N - x^* \Vert^2 &\le \frac{\Vert x_0 -x^* \Vert^2}{2^N} + \frac{2 \sigma^2}{\alpha \beta}. \label{eq:EpochSGD-xN}
\end{align}
For the subsequent $K$ epochs, by the parameter setting in Algorithm \ref{alg:SGD-SC} and \Eqref{eq:each-epoch}, we have that 
\begin{align*}
   \BE \Vert x_{N+k} - x^* \Vert^2 \le \frac{\Vert x_{N+k-1} - x^* \Vert^2}{4} + \frac{\sigma^2}{2^{k-1} \alpha \beta}.
\end{align*}
Then it is easy to prove by induction that
\begin{align} \label{eq:induction-SGD}
    \BE \Vert x_{N+k} - x^* \Vert^2 \le \frac{\Vert x_0 - x^* \Vert^2}{2^{N+2k}} + \frac{\sigma^2}{2^{k-2} \alpha \beta}.
\end{align}
In the above, the induction base follows \Eqref{eq:EpochSGD-xN}. If we suppose that \Eqref{eq:induction} holds for $k=j-1$, then for $k=j$ we have that
\begin{align*}
    \BE \Vert x_{N+j+1} - x^* \Vert^2 &\le \frac{1}{4} \cdot \left( \frac{\Vert x_0 -x^* \Vert^2}{2^{N+2(j-1)}} + \frac{\sigma^2}{2^{j-3} \alpha \beta} \right) + \frac{\sigma^2}{2^{j-1} \alpha \beta} \\
    &\le \frac{\Vert x_0 - x^* \Vert^2}{2^{N+2j}} + \frac{\sigma^2}{2^{j-2} \alpha \beta},
\end{align*}
which completes the induction of \Eqref{eq:induction-SGD}.
\end{proof}
}


% Below, we integrate the momentum technique in the outer loop:
% \begin{align*} 
%     g_{t+1} = (1- \gamma) g_t + \gamma \hat \nabla \fL_{\lambda}^*(x_t)
% \end{align*}
% }
% Define the variance term as 
% \begin{align*}
%     \Delta g_t : =  g_t - \nabla_x \fL_{\lambda}(x_t, y_t^K, z_t^K).
% \end{align*}
% Then we have that
% \begin{align*}
%     \Delta g_{t+1} & =  g_{t+1} - \nabla \fL_{\lambda}^*(x_{t+1})  \\
%     &= (1- \gamma) g_t + \gamma \hat \nabla \fL_{\lambda}^*(x_t)  - \nabla \fL_{\lambda}^*(x_{t+1})  \\
%     &= (1- \gamma) g_t + \gamma \nabla \fL_{\lambda}^*(x_t)  - \nabla \fL_{\lambda}^*(x_{t+1}) + \gamma ( \hat \nabla \fL_{\lambda}^*(x_t) - \nabla \fL_{\lambda}^*(x_t)) \\
%     &= (1- \gamma) \Delta g_t + (\nabla \fL_{\lambda}^*(x_t) - \nabla \fL_{\lambda}^*(x_{t+1}))  + \gamma ( \hat \nabla \fL_{\lambda}^*(x_t) - \nabla \fL_{\lambda}^*(x_t))
% \end{align*}
% Taking square on both sides leads to
% \begin{align*}
%     \Vert \Delta g_{t+1} \Vert^2 &= (1 - \gamma)^2 \Vert \Delta g_t \Vert^2 + \Vert  \nabla \fL_{\lambda}^*(x_t) - \nabla \fL_{\lambda}^*(x_{t+1})  + \gamma ( \hat \nabla \fL_{\lambda}^*(x_t) - \nabla \fL_{\lambda}^*(x_t)) \Vert^2 \\
%     &\quad + 2 (1- \gamma) \Delta g_t^\top  ( \nabla \fL_{\lambda}^*(x_t) - \nabla \fL_{\lambda}^*(x_{t+1}) )+ 2 \gamma (1- \gamma) \Delta g_t^\top ( \hat \nabla \fL_{\lambda}^*(x_t) - \nabla \fL_{\lambda}^*(x_t)).
% \end{align*}
% Applying Young's inequality, we further have that
% \begin{align*}
%     \Vert \Delta g_{t+1} \Vert^2 &\le  (1 - \gamma)^2 \Vert \Delta g_t \Vert^2 + 2 \Vert \nabla \fL_{\lambda}^*(x_t) - \nabla \fL_{\lambda}^*(x_{t+1}) \Vert^2 + 2 \gamma^2 \Vert \hat \nabla \fL_{\lambda}^*(x_t) - \nabla \fL_{\lambda}^*(x_t) \Vert^2 \\
%     &\quad + \frac{\gamma}{2} \Vert \Delta g_t \Vert^2 + \frac{2}{\gamma} \Vert \nabla \fL_{\lambda}^*(x_t) - \nabla \fL_{\lambda}^*(x_{t+1}) \Vert^2 + 2 \gamma (1- \gamma) \Delta g_t^\top ( \hat \nabla \fL_{\lambda}^*(x_t) - \nabla \fL_{\lambda}^*(x_t)).
% \end{align*}
% Taking expectation and using the $L$-smoothness of $ \fL_{\lambda}^*(x)$ yields that
% \begin{align*}
%     \BE \Vert \Delta g_{t+1} \Vert^2 &\le (1- \gamma) \Vert \Delta g_t \Vert^2 +  2 \gamma^2 \BE \Vert \hat \nabla \fL_{\lambda}^*(x_t) - \nabla \fL_{\lambda}^*(x_t) \Vert^2 \\
%     &\quad + \left( \gamma^2 - \frac{\gamma}{2} \right) \Vert \Delta g_t \Vert^2 + \frac{4 L^2}{\gamma} \Vert x_{t+1} - x_t \Vert^2 + 2 \gamma (1- \gamma) \Delta g_t^\top ( \BE \hat \nabla \fL_{\lambda}^*(x_t) - \nabla \fL_{\lambda}^*(x_t)).
% \end{align*}
% Recalling the descent lemma
% \begin{align*} 
%    \BE[ \fL_{\lambda}^*(x_{t+1})]  \le \BE \left[\fL_{\lambda}^*(x_t) 
%     - \frac{\eta_x}{2} \Vert \nabla \fL_{\lambda}^*(x_t) \Vert^2 - 
%     \frac{1}{4\eta_x} \Vert x_{t+1} - x_{t} \Vert^2
%     + \frac{\eta_x}{2} \Vert  \Delta g_t \Vert^2 \right].
% \end{align*}
% Next, we show that the logarithmic factor in $\epsilon$ can be removed by simply
% adding a normalization step in $x$. Below, we recall the descent lemma for normalized gradient descent.

% \begin{lem}[Lemma 2 by \citet{cutkosky2020momentum}] Suppose $h(x) : \BR^d \rightarrow \BR$ is $\beta$-gradient Lipschitz. For any given stepsize $\eta>0$, consider the following update of normalized gradient descent:
% \begin{align*}
%     x_{t+1} = x_t - \eta \frac{g_t}{\Vert g_t \Vert}.
% \end{align*}
% Then it holds that
% \begin{align*}
%     h(x_{t+1}) \le h(x_t) - \frac{\eta}{3}  \Vert \nabla h(x_t) \Vert + \frac{8 \eta}{3} \Vert g_t - \nabla h(x_t) \Vert^2 + \frac{L \eta^2}{2}.
% \end{align*}
% \end{lem}


\section{Proofs of Finding Second-Order Stationarity} \label{apx:2nd}

% In this section, we focus on how to escape saddle points in bilevel optimization. 
The following theorem generalizes the result of Perturbed Gradient Descent~\citep{jin2017escape}
to Inexact Perturbed Gradient Descent by
allowing an inexact gradient $\hat \nabla h(x)$ that is close to the exact gradient $\nabla h(x)$.

\begin{algorithm*}[htbp]  
\caption{Inexact Perturbed Gradient Descent} \label{alg:iPGD}
\begin{algorithmic}[1] 
\STATE \textbf{for} $t = 0,1,\cdots,T-1$ \\
\STATE \quad \textbf{if} $ \Vert \hat \nabla h(x_t) \Vert \le \frac{4}{5} \epsilon$ \textbf{and} no perturbation added in the last $\fT$ steps \\[1mm]
\STATE \quad \quad $x_{t} = x_t - \eta \xi_t $, where $\xi_t \sim \sB(r)$ \\[1mm]
\STATE \quad \textbf{end if} \\[1mm]
\STATE \quad $x_{t+1} = x_t -  \eta  \hat \nabla h(x_t)$ \\[1mm]
\STATE \textbf{end for} \\[1mm]
\end{algorithmic}
\end{algorithm*}

\begin{thm}[{\citet[Lemma 2.9]{huang2022efficiently}}]
    \label{thm:iPGD}
Suppose $h(z): \BR^d \rightarrow \BR$ is
$L$-gradient Lipschitz and $\rho$-Hessian Lipschitz. 
Set the parameters in Algorithm \ref{alg:iPGD} as 
\begin{align} \label{cond:iota}
    \eta = \frac{1}{L}, \quad r = \frac{\epsilon}{400 \iota^3}, \quad \fT = \frac{L}{\sqrt{\rho \epsilon}} \cdot \iota ~~{\rm with }~~ \iota \ge 1 ~~ {\rm and} ~~ \delta \ge  \frac{L \sqrt{d}}{\sqrt{\rho \epsilon}} \iota^2 2^{8-\iota}.
\end{align}
Once the following condition holds in each
iteration:
\begin{align} \label{cond:PF2BA}
    \Vert \hat \nabla h(x_t) - \nabla h(x_t) \Vert \le \underbrace{\min \left\{ 
    \frac{1}{20 \iota^2},  \frac{1}{16 \iota^2 2^{\iota}} \right\} \epsilon}_{:=\zeta},
\end{align}
then we can find a point $x_t$ satisfying
\begin{align*}
    \Vert \nabla h(x_t) \Vert \le \epsilon, \quad  \nabla^2 h(x_t)  \succeq - \sqrt{\rho \epsilon}~ I_{d}
\end{align*}
within $T= \fO \left( {\iota^4  \Delta L}{\epsilon^{-2}}   \right)$ iterations with probability $1-\delta$, where
$\Delta = h(x_0) - \inf_{x \in \BR^{d}} h(x)$.    
\end{thm}

Then we can show the convergence of perturbed F${}^2$BA.
\thmPFBA*
\begin{proof}
% The result follows from
% Lemma \ref{lem:1st}, Lemma \ref{lem:2rd}, Proposition \ref{prop:iPGD} and Proposition \ref{prop:error}.
Let $L$ be the gradient Lipschitz coefficient and $\rho$ be the Hessian Lipschitz coefficient of $\fL_{\lambda}^*(x)$, respectively.
According to Lemma \ref{lem:1st}, Lemma \ref{lem:2rd}, Lemma \ref{lem:yy-lambda}  and the setting of $\lambda$, we have
\begin{enumerate}[label=\alph*.]
    \item $\sup_{x \in \BR^{d_x}} \Vert \nabla \fL_{\lambda}^*(x) - \nabla \varphi(x) \Vert = \fO(\epsilon)$. 
    \item $\sup_{x \in \BR^{d_x}} \Vert \nabla^2 \fL_{\lambda}^*(x) - \nabla^2 \varphi(x) \Vert = \fO(\sqrt{\rho \epsilon})$.
    \item $\fL_{\lambda}^*(x_0) - \inf_{x \in \BR^{d_x}}  \fL_{\lambda}^*(x) = \fO(\Delta)$.
    \item $\Vert y_0 - y_{\lambda}^*(x_0) \Vert^2 + \Vert y_0 - y^*(x_0) \Vert^2 =\fO(R)$.
    \item $ L:=\sup_{x \in \BR^{d_x}}\Vert \nabla^2 \fL_{\lambda}^*(x) \Vert = \fO(\ell \kappa^3)$.
    \item $\rho: = \sup_{x \in \BR^{d_x}} \Vert \nabla^3 \fL_{\lambda}^*(x) \Vert = \fO(\ell \kappa^5) $.
\end{enumerate}
Then it suffices to show that the algorithm can find an $\epsilon$-second-order stationary point of $\fL_{\lambda}^*(x)$ under our choice of parameters.
We apply Theorem \ref{thm:iPGD} to prove this result. 
Recall $\zeta$ defined in \Eqref{cond:PF2BA}.
It remains to show that 
\begin{align} \label{eq:cond}
    \Vert \nabla \fL_{\lambda}^*(x_t) - \nabla \fL_{\lambda}^*(x_t) \Vert \le \zeta
\end{align}
holds for all $t$. Recall that
\begin{align*}
    \Vert \hat \nabla \fL_{\lambda}^*(x_t) - \nabla \fL_{\lambda}^*(x_t) \Vert \le 2 \lambda L_g \Vert y_{t}^K- y_{\lambda}^*(x_t) \Vert + \lambda L_g \Vert z_{t}^K - y^*(x_t) \Vert.
\end{align*} 
 If we let
\begin{align} \label{eq:Kt}
    K_t = \frac{4 L_g}{\mu} \log \left( 
    \frac{4 \lambda L_g \delta_t}{\zeta} 
    \right), 
\end{align}
where $\delta_t \ge \Vert y_t^0 - y_{\lambda}^*(x_t) \Vert^2 + \Vert z_t^0 - y^*(x_t) \Vert^2  $.
Then by Theorem \ref{thm:GD-SC}, we have
\begin{align} \label{eq:induction}
     \max \left\{ \Vert y_t^K - y_{\lambda}^*(x_t) \Vert , ~\Vert z_t^K-  y^*(x_t) \Vert   \right\} \le \frac{\zeta}{4 \lambda L_g},
\end{align}
which implies \Eqref{eq:cond}. 
%The only remaining issue is that the algorithm does not know the value of  in \Eqref{eq:Kt}. We address this issue by showing a upper bound of it.
% Suppose that we have already upper bounded $\delta_t$, we can then use it to upper bound $\delta_{t+1}$.
By \Eqref{eq:delta_t}, if we let $K_t \ge 8L_g/\mu$ for all $t$, then we can show that
\begin{align}  \label{eq:dekta-recur}
     \delta_{t+1}
    &\le \frac{1}{2} \delta_t +\frac{34 L_g^2}{\mu^2} \Vert x_{t+1} - x_t \Vert^2.
\end{align}
% Note that the value of $\Vert x_{t+1} - x_t \Vert^2$ is known in advance. This recursion gives a valid upper bound of $\delta_t$ that can be used to specify $K_t$ in the algorithm. 
% Algorithm using this choice of $K_t$ can provably find an $\epsilon$-second-order stationary point of $\fL_{\lambda}^*(x)$ within 
The total number of iterations can be bounded by
\begin{align} \label{plug:2}
\sum_{t=0}^{T-1} K_t \le \frac{4L_g}{\mu} \sum_{t=0}^{T-1} \log \left( 
\frac{4 \lambda L_g \delta_t}{\zeta}
\right)   \le \frac{4L_g T}{\mu} \log \left( \frac{4 \lambda L_g \sum_{t=0}^{T-1}\delta_t}{\zeta T} \right)
\end{align}
 first-order oracle calls. It remains to show that $\sum_{t=0}^{T-1} \delta_t$ is bounded.
Telescoping over $t$ in \Eqref{eq:dekta-recur}, we obtain
\begin{align} \label{plug:1}
    \sum_{t=0}^{T-1} \delta_{t} \le 2 \delta_0 + \frac{68L_g^2}{\mu^2} \sum_{t=0}^{T-1} \Vert x_{t+1} - x_t \Vert^2.
\end{align}
By \Eqref{eq:final}, if we let $K_t \ge \Omega( \kappa \log(\lambda \kappa))$ for all $t$, we can obtain
\begin{align*}
\frac{1}{8 \eta} \sum_{t=0}^{T-1} \Vert x_{t+1} - x_t \Vert^2 \le \fL_{\lambda}^*(x_0) - \inf_{x \in \BR^{d_x}} \fL_{\lambda}^*(x) + \Vert y_0 - y_{\lambda}^*(x_0) \Vert^2 + \Vert y_0 - y^*(x_0) \Vert^2 .
\end{align*}
Plugging into \Eqref{plug:1} and then \Eqref{plug:2} yields the upper complexity bound of $\tilde \fO(\ell \kappa^4 \epsilon^{-2})$ as claimed. 
\end{proof}

The following theorem generalizes the result of Nonconvex AGD~\citep{li2023restarted} by allowing an inexact gradient. To simplify the description, we define one round of the algorithm between
two successive restarts to be one ``epoch'' by following \citet{li2023restarted}.

\begin{algorithm*}[htbp]  
\caption{Inexact Nonconvex Accelerated Gradient Descent} \label{alg:iAGD}
\begin{algorithmic}[1] 
\STATE $x_{-1} = x_0 $ \\
\STATE \textbf{while} $t <T$ \\
\STATE \quad $x_{t+1/2} = x_t + (1- \theta)(x_t - x_{t-1})$ \\
\STATE \quad $x_{t+1} = x_{t+1/2} - \eta \hat \nabla h(x_{t+1/2}) $ \\
\STATE \quad $t = t+1$ \\
\STATE \quad \textbf{if} $t \sum_{j=0}^{t-1} \Vert x_{j+1} - x_j \Vert^2 > B^2$ \\
\STATE \quad \quad $t=0$, $x_{-1} = x_0 = x_t + \xi_t \vone_{ \Vert \hat \nabla h(x_{t+1/2}) \Vert \le \frac{B}{2 \eta} }, ~{\rm where} ~~\xi_t \sim \sB(r)$ \\
\STATE \quad \textbf{end if} \\
\STATE \textbf{end while} \\
\STATE $T_0 = \arg \min_{ \lfloor \frac{T}{2} \rfloor \le t \le T-1} \Vert x_{t+1} - x_t \Vert$ \\
\STATE \textbf{return} $ x_{\rm out} = \frac{1}{T_0+ 1} \sum_{t=0}^{T_0} x_{t+1/2}$ \\
\end{algorithmic}
\end{algorithm*}



\begin{thm}[{\citet[Theorem 4.1]{yang2023accelerating}}]
    \label{thm:iAGD}
Suppose $h(z): \BR^d \rightarrow \BR$ is
$L$-gradient Lipschitz and $\rho$-Hessian Lipschitz. Assume that $\sqrt{\epsilon \rho } \le L$.
Set the parameters in Algorithm \ref{alg:iAGD} as 
\begin{align*}
    &\chi = \fO( \log (\nicefrac{d}{\delta \epsilon})),~~ \eta = \frac{1}{4 L}, ~~ T = \frac{2 \chi}{\theta}, ~ \theta = \frac{1}{2} ( \rho \epsilon \eta^2 )^{1/4}<1,\\
    & B = \frac{1}{288 \chi^2} \sqrt{\frac{\epsilon}{\rho}}, ~~ 
     r =\min \left\{ \frac{B + B^2}{\sqrt{2}}, \frac{\theta B}{20 T}, \frac{\theta B^2}{2T} \right\} = \fO(\epsilon).
\end{align*}
% \begin{align} \label{cond:iota}
%     \eta = \frac{1}{L}, \quad r = \frac{\epsilon}{400 \iota^3}, \quad \fT = \frac{L}{\sqrt{\rho \epsilon}} \cdot \iota ~~{\rm with }~~ \iota \ge 1 ~~ {\rm and} ~~ \delta \ge  \frac{L \sqrt{d}}{\sqrt{\rho \epsilon}} \iota^2 2^{8-\iota}.
% \end{align}
Once the following condition holds in each iteration:
\begin{align} \label{cond:RAGD}
     \Vert \hat \nabla h(x_t) - \nabla h(x_t) \Vert \le \underbrace{\min \left\{ \frac{\rho B \zeta r \theta}{2 \sqrt{d}}, \epsilon^2 \right\}}_{:=\zeta}.
\end{align}
% \begin{align*} \label{cond:PF2BA}
%     \Vert \hat \nabla h(x_t) - \nabla h(x_t) \Vert \le \underbrace{\min \left\{ 
%     \frac{1}{20 \iota^2},  \frac{1}{16 \iota^2 2^{\iota}} \right\} \epsilon}_{:=\zeta},
% \end{align}
The algorithm ensures that
\begin{align} \label{eq:des-AGD}
    h(x_{\fT}) - h(x_0) \le - \frac{\epsilon^{1.5}}{663552 \sqrt{\rho} \chi^5},
\end{align}
where $\fT$ denotes the iteration number when ``if'' condition in Algorithm \ref{alg:iAGD} (Line 6) triggers:
\begin{align*}
\fT = \min_{t'} \left\{ t' \mid t' \sum_{t=0}^{t'-1} \Vert x_{t+1} - x_t \Vert^2 > B^2 \right\}.
\end{align*}
Therefore, the algorithm terminates in at most $\fO(\Delta \sqrt{\rho} \chi^5 \epsilon^{-3/2})$ epochs and $ \fO( \Delta \sqrt{L} \rho^{1/4} \chi^6 \epsilon^{-1.75})$ iterations. In addition, the output satisfies
\begin{align*}
    \Vert \nabla h(x_t) \Vert \le \epsilon, \quad  \nabla^2 h(x_t)  \succeq - \sqrt{\rho \epsilon}~ I_{d}, 
\end{align*}
with probability $1-\delta$, where
$\Delta = h(x_0) - \inf_{x \in \BR^{d}} h(x)$.    
\end{thm}

\thmAccFBA*

\begin{proof}
Let $T'$ be the number of total iterations. We renumber the iterates in the algorithm as $\{x_{t} \}_{t=0}^{T'-1}$,  $\{x_{t+1/2} \}_{t=0}^{T'-1}$, and let the corresponding inner loop number as $\{K_t \}_{t=0}^{T'-1}$.
Identical to the proof of Theorem \ref{thm:PF2BA} (with a different definition of $\zeta$ that only effects the logarithmic factors),
we can set $K_t = \fO(\kappa \log (\nicefrac{\lambda L_g \delta_t}{\zeta}))$ to
ensure Condition \ref{cond:RAGD} to hold, where 
$\delta_t \ge \Vert y_t^0 - y_{\lambda}^*(x_{t+1/2}) \Vert^2 + \Vert z_t^0 - y^*(x_{t+1/2}) \Vert^2 $ is ensured recursively identical to \Eqref{eq:dekta-recur}.
By \Eqref{plug:2} we have
\begin{align*}
    \sum_{t=0}^{T'-1} K_t \le \frac{4L_g T}{\mu} \log \left( \frac{4 \lambda L_g \sum_{t=0}^{T'-1}\delta_t}{\zeta T'} \right)
\end{align*}
It remains to show that $\sum_{t=0}^{T'-1} \delta_t$ is bounded.  By \Eqref{plug:1}, it suffices to show that $\sum_{t=0}^{T'-1} \Vert x_{t+3/2} - x_{t+1/2} \Vert^2  $ is bounded. Since $x_{t+1/2} = x_t + (1-\theta) (x_t - x_{t-1})$, it suffices to show that $\sum_{t=0}^{T'-1} \Vert x_{t+1} - x_t \Vert^2$ is bounded.
Because the number of epochs is finite by Theorem \ref{thm:iAGD}, it suffices to bound the term in each epoch. 
We know that
\begin{align*}
    &\quad \sum_{t=0}^{\fT -1} \Vert x_{t+1} - x_t \Vert^2  \\
    &= \sum_{t=0}^{\fT -2} \Vert x_{t+1} - x_t \Vert^2  + \Vert x_{\fT} - x_{\fT-1} \Vert^2 \\
    &\le B^2 + \Vert x_{\fT} - x_{\fT-1} \Vert^2.
\end{align*}
We continue to bound  the last term $\Vert x_{\fT} - x_{\fT-1} \Vert^2$.
Since we have
\begin{align*}
    x_{t+1} - x_{t} &= (1- \theta) (x_t - x_{t-1}) - \eta_x \hat \nabla \fL_{\lambda}^*(x_{t+1/2}).
\end{align*}
It remains to upper bound $ \Vert \hat \nabla \fL_{\lambda}^*(x_{\fT-1/2}) \Vert$. 
As $ \Vert \hat \nabla \fL_{\lambda}^*(x_{\fT-1/2}) - \nabla \fL_{\lambda}^*(x_{\fT-1/2}) \Vert = \fO(\zeta)$, it remains to upper bound $\Vert \nabla \fL_{\lambda}^*(x_{\fT-1/2}) \Vert$. 
By (18) in \citet{yang2023accelerating}, we have that
\begin{align*}
    \fL_{\lambda}^*(x_{\fT}) \le \fL_{\lambda}^*(x_0) + \frac{B^2}{8 \eta_x} - \frac{3 \eta_x}{8} \Vert \nabla \fL_{\lambda}^*(x_{\fT-1/2}) \Vert^2 + B \zeta + \frac{5 \eta_x \zeta^2 \fT}{8}.
\end{align*}
This implies that $\Vert \nabla \fL_{\lambda}^*(x_{\fT-1/2}) \Vert$ is bounded once $\Delta$ is bounded.
\end{proof}

\section{Implement Our Algorithms in the Distributed Setting} \label{apx:dist}

Our algorithms can also be easily applied to the distributed scenario, when both the upper and lower-level functions adopt the following finite-sum structure:
\begin{align}  \label{prob:fn}
    f(x,y) := \frac{1}{m} \sum_{i=1}^m f_i(x,y), \quad g(x,y) := \frac{1}{m} \sum_{i=1}^m g_i(x,y),
\end{align}
where $f_i$ and $g_i$ denote the local function on the $i$-th agent. 
% One of the primal challenges for HVP-based methods in distributed bilevel optimization arises from the Hessian inverse in the expression of hyper-gradient (\Eqref{hyper-grad}) since directly aggregating the Hessian matrix requires an expensive $\fO(d_{y}^2)$ communication complexity~\citep{yangdecentralized,lu2022stochastic}. To address this issue, existing works reformulate the Hessian inverse as a sub-problem and solve it with distributed algorithms that only require communicating HVPs~\citep{tarzanagh2022fednest,chen2022decentralized,xiao2023communication}.
% In contrast, the extension of F${}^2$BA to the distributed setting (Algorithm~\ref{alg:CF2BA}) naturally avoids this additional distributed sub-solver for matrix inverse. 
% \subsection{Distributed F${}^2$BA}
% As we have discussed, gradient-based methods have superiority over HVP-based methods when the HVP oracle is expensive.
% One of these scenarios is distributed bilevel optimization, where both $f$ and $g$ are distributed on $m$ agents as \Eqref{prob:fn}. The direct extension of HVP-based methods is to aggregate the local hyper-gradients in \Eqref{hyper-grad} on all the agents. Unfortunately, this simple algorithm may not converge and additional techniques are required to address this challenge~\citep{xiao2023communication,tarzanagh2022fednest,chen2022decentralized}. In contrast, the gradient-based algorithm F${}^2$BA does not have this issue since 
% aggregating the local gradients in \Eqref{eq:nabla-Lag} gives the exact 
% form of $\nabla \fL_{\lambda}^*(x)$.
% Before presenting our algorithm, we first
% introduce the notations for the distributed F${}^2$BA. 
Under this setting, each agent $i$ has its own local variables within the optimization algorithm: $X(i) \in \BR^{d_x}, Y(i) \in \BR^{d_y},Z(i) \in \BR^{d_y}$.
For the convenience of presenting the algorithm, we aggregate all the local variables (denoted by row vectors) in one matrix and denote
\begin{align*}
X=\begin{bmatrix}
X(1) \\ \vdots \\ X(m)
\end{bmatrix}&\in\BR^{m\times d_x},
\quad
Y=\begin{bmatrix}
Y(1) \\ \vdots \\ Y(m)
\end{bmatrix}\in\BR^{m\times d_y}
\quad \text{and} \quad
Z=\begin{bmatrix}
Z(1) \\ \vdots \\ Z(m)
\end{bmatrix}\in\BR^{m\times d_y},
\end{align*}
We denote~$\vone = [1,\cdots,1]^\top \in \BR^m$ and use the lowercase with the bar to represent the mean vector, such as $\bar x = \frac{1}{m} \vone^\top X$.
Let $d = d_x+d_y$ and we similarly define the aggregated gradients 
{\small \begin{align*}
\nabla \vf(X,Y)=\begin{bmatrix}
\nabla f_1(X(1),Y(1)) \\ \vdots \\ \nabla f_m(X(m),Y(m))
\end{bmatrix}
\in\BR^{m\times d}, ~\nabla \vg(X,Y)=\begin{bmatrix}
\nabla g_1(X(1),Y(1)) \\ \vdots \\ \nabla g_m(X(m),Y(m))
\end{bmatrix}
\in\BR^{m\times d}.    
\end{align*}}


\begin{algorithm*}[t]  
\caption{Distributed F${}^2$BA $\left(\bar x_0,\bar y_0, \eta_x,\eta_y,\eta_z, \lambda, T,K\right)$} \label{alg:CF2BA}
\begin{algorithmic}[1] 
\STATE $ X_0 = \vone  \bar x_0, ~ Y_0 = \vone  \bar y_0,~ Z_0 = \vone  \bar y_0$ \\[1mm]
\STATE \textbf{for} $ t =0,1,\cdots,T-1 $ \\[1mm]
\STATE \quad $ Y_t^0 =  Y_{t}, ~ Z_t^0(i) = Z_{t}$ \\[1mm]
\STATE \quad \textbf{for} $ k =0,1,\cdots,K-1$ \\[1mm]
\STATE \quad \quad  $ V_t^{k} = \lambda \nabla_y \vg(X_t, Z_t^k)$, \quad  $U_t^{k} =  \nabla_y \vf(X_t,Y_t^k) +  \lambda  \nabla_y \vg(X_t,Y_t^k) $ \\[1mm]
\STATE \quad \quad Aggregate and broadcast $\bar v_t^{k} = \frac{1}{m} \vone \vone^\top V_t^k, $~ $\bar u_t^{k} = \frac{1}{m} \vone   \vone^\top U_t^k$ \\[1mm]
\STATE \quad \quad  $ Z_t^{k+1} = Z_t^k - \eta_z \vone   \bar v_t^k   $ ,\quad  $ Y_t^{k+1} = Y_t^k - \eta_y \vone \bar u_t^k   $ \\[1mm]
\STATE \quad \textbf{end for} \\[1mm]
% \STATE \quad $Z_{t+1} = Z_t^K,~ Y_{t+1} = Y_t^K $ \\[1mm]
\STATE \quad $ H_t= \nabla_x \vf(X_t,Y_{t}^K) + \lambda ( \nabla_x \vg(X_t,Y_{t}^K) - \nabla_x \vg(X_t,Z_{t}^K) )$ \\[1mm]
\STATE \quad Aggregate and broadcast $\bar h_t = \frac{1}{m} \vone   \vone^\top H_t$ \\[1mm]
\STATE \quad $X_{t+1} = X_t -  \eta_x  \vone  \bar h_t$ \\[1mm]
\STATE \textbf{end for} \\[1mm]
\end{algorithmic}
\end{algorithm*}


% Some applications of bilevel optimization like neural architecture search suffer from a heavy computational burden, so the experiments are usually conducted in a distributed way~\citep{chen2019progressive}.
% In this section, we extend
% F${}^2$BA to the distributed setting, where both the upper and lower-level functions are distributed on $m$ agents as \Eqref{prob:fn}. The detailed procedure is presented in Algorithm \ref{alg:CF2BA}.

The classic paradigm for distributed algorithms is to compute the local gradients on each agent in parallel, and then aggregate them on the server. 
However, challenges arise when extending existing HVP-based methods for bilevel optimization from the single-machine setting to the distributed setting. Given a variable $\bar x$, one may want to calculate its hyper-gradient (\Eqref{hyper-grad}) according to this paradigm by:
\begin{align*}
    \hat \nabla \varphi(\bar x) = {\bf \Pi} \left( \nabla_x \vf(\vone x, \vone y^*(\bar x)) - \nabla_{xy}^2 \vf(\vone x,\vone y^*(\bar x)) [\nabla_{yy}^2 \vg(\vone \bar x, \vone y^*(\bar x))]^{-1} \nabla_y \vf(\vone \bar x, \vone y^*(\bar x))  \right),
\end{align*}
where $y^*(x) := \arg \min_{y \in \BR^{d_y}} g(x,y)$ and ${\bf \Pi} := \frac{1}{m} \vone \vone^\top$ denotes the aggregation operator on the server. 
But the nested structure of the hyper-gradient indicates that $ \hat \nabla \varphi(\bar x) \ne \nabla \varphi(\bar x)$. As a consequence, 
researchers need to pay extra effort to make HVP-based methods work in distributed bilevel problems~\citep{tarzanagh2022fednest,chen2022decentralized}.
It is a non-trivial issue to address, especially when the operator $ {\bf \Pi} \nabla_{yy}^2 \vg(\vone \bar x,\vone y^*(\bar x)) $ is forbidden to avoid an unacceptable $\fO(d_y^2)$ communication complexity.

In contrast, the distributed extension of F${}^2$BA naturally avoid this challenge since
\begin{align*}
    \nabla \fL_{\lambda}^*(\bar x) = {\bf \Pi} \left( \nabla_x \vf(\vone x,\vone y_{\lambda}^*(\bar x)) + \lambda (\nabla_x \vg(\vone \bar x, \vone y_{\lambda}^*(\bar x) ) - \nabla_x \vg (\vone \bar x, \vone y^*(\bar x)) ) \right),
\end{align*}
where $y_{\lambda}^*(x) := \arg \min_{y \in \BR^{d_y}} f(x,y) + \lambda g(x,y)$. It means that the previously mentioned classic aggregation paradigm for distributed optimization directly works for F${}^2$BA without any additional modification in the algorithm, as stated below.

% The most natural idea is to aggregate the local gradients in \Eqref{eq:nabla-Lag} on all agents. This procedure is presented in 
% Algorithm \ref{alg:CF2BA}.
% Thanks to the form of \Eqref{eq:nabla-Lag}, one can easily verify that the mean variables in Algorithm \ref{alg:CF2BA} follow the same trajectories as in Algorithm \ref{alg:F2BA}.
\begin{prop} 
Running Algorithm \ref{alg:CF2BA} is equivalent to running  Algorithm \ref{alg:F2BA}  on the mean variables
\begin{align*}
    \bar x = \frac{1}{m} \vone^\top X,~ \bar y = \frac{1}{m} \vone^\top Y ~{\rm and}~ \bar z = \frac{1}{m} \vone^\top Z.
\end{align*}
\end{prop}

Then the convergence of Algorithm \ref{alg:CF2BA} directly follows Theorem \ref{thm:F2BA}.
%that for F${}^2$BA distributed bilevel problems are as easy as the single-machine setup.  As we have discussed, such an advantage does not hold for second-order methods.
\begin{cor} \label{cor:CF2BA}
Suppose Assumption \ref{asm:sc} holds.
Algorithm \ref{alg:CF2BA} with the same parameters in Theorem \ref{thm:F2BA} can find  $X_t$ satisfying $ \Vert \nabla \varphi(\bar x_t) \Vert \le \epsilon $ within $\fO(\ell \kappa^4 \epsilon^{-2} \log(\nicefrac{\ell \kappa}{\epsilon})) $ {iterations} and  $\fO( \ell \kappa^4 \epsilon^{-2} \log(\nicefrac{\ell \kappa}{\epsilon})) $ communication rounds.
\end{cor}

Algorithm \ref{alg:CF2BA} is a near-optimal distributed algorithm since both the iteration and communication complexity match the $\Omega(\epsilon^{-2})$ lower bound (Theorem 1 by \citet{lu2021optimal}), up to logarithmic factors.
Compared with the HVP-based methods, the distributed F${}^2$BA is more
practical since it
neither requires a $\fO(d_y^2)$ communication complexity per iteration~\citep{yangdecentralized} nor an additional distributed sub-solver for matrix inverse~\citep{tarzanagh2022fednest,chen2022decentralized}. 

{We remark that all algorithms in our main text can be implemented in a similar way as we implement F${}^2$BA in the distributed scenario.
}

\section{Discussions on Closely Related Works}

In this section, we discuss our method with closely related works, including \citep{shen2023penalty,kwon2023fully}.
Both our work and theirs use a penalty method. 
The main difference in algorithm design is that their method uses a single-time-scale update in $(x,y)$ while we use a two-time-scale update to improve their $\fO( \lambda \epsilon^{-2} \log(\nicefrac{1}{\epsilon}))$ first-order complexity to $\fO(\epsilon^{-2} \log \lambda )$, where $\lambda \asymp \epsilon^{-1}$ due to Lemma \ref{lem:1st}.

In the following, we give a more detailed comparison of another aspects, including the optimality notion and assumptions of our work and theirs.

\subsection{Discussions on the Weaker Notions in \citep{shen2023penalty}} \label{apx:diss-shen}

\citet{shen2023penalty} also proposed a similar penalty-based algorithm, and proved that the penalty method can find an $\epsilon$-stationary point (formally defined below) of $\fL_{\lambda}(x,y)$ in $\fO(\lambda \epsilon^{-2} \log (\nicefrac{1}{\epsilon}))$ first-order oracle calls.
Below, we show that their result implies a $\fO(\epsilon^{-3} \log (\nicefrac{1}{\epsilon})$ first-order complexity to find an $\epsilon$-stationary point of $\varphi(x)$. 


\begin{dfn} \label{dfn:sta-Lambda}
We say $(x,y)$ is an $\epsilon$-stationary point of $\fL_{\lambda}(x,y)$ if $\Vert \nabla \fL_{\lambda}(x,y) \Vert \le \epsilon$.
\end{dfn}

We show that their notion and ours can be translated in both directions.

\begin{lem} \label{lem:translate}
Suppose Assumption \ref{asm:sc} and \ref{asm:stochastic} hold. Set $\lambda$ as Theorem \ref{thm:F2BSA}.
\begin{enumerate}
    \item If a point $x$ is an $\epsilon$-stationary point of $\varphi(x)$ in terms of Definition \ref{dfn:sta-hyper}, then we can find a point $y$ such that $(x,y)$ is an $2\epsilon$-stationary point of $\fL_{\lambda}(x,y)$ in terms of Definition \ref{dfn:sta-Lambda} with $\fO(\kappa \log(\nicefrac{\kappa}{\epsilon}))$,
    $\fO(\kappa^2 \epsilon^{-2})$, or $\fO(\kappa^8 \epsilon^{-2})$ first-order oracle calls under the deterministic $(\sigma_f = \sigma_g=0)$, partially stochastic ($\sigma_f >0, \sigma_g=0$) and fully stochastic setting $(\sigma_f >0, \sigma_g>0)$ respectively.
    \item If a point $(x,y)$ is an $\epsilon$-stationary point of $\fL_{\lambda}(x,y)$ in terms of Definition \ref{dfn:sta-Lambda}, then it is a $\fO(\kappa \epsilon)$-stationary point of $\varphi(x)$ in terms of Definition \ref{dfn:sta-hyper}.
\end{enumerate}
\end{lem}

\begin{proof}
Note that $\varphi(x) = \min_{y \in \BR^{d_y}} \fL_{\lambda}(x,y)$. Then by Danskin's lemma, we have $\nabla \varphi(x) = \nabla_x \fL_{\lambda}(x,y_{\lambda}^*(x))$. Therefore, we further have
\begin{align*}
    \Vert \nabla \varphi(x) \Vert &\le \Vert \nabla_x \fL_{\lambda}(x,y) \Vert + 2 \lambda L_g \Vert y_{\lambda}^*(x) - y^*(x) \Vert \\
    &\le \Vert \nabla_x \fL_{\lambda}(x,y) \Vert + \frac{4 L_g}{\mu} \Vert \nabla_y \fL_{\lambda}(x,y) \Vert,
\end{align*}
where we use the fact that $\fL_{\lambda}(x,y)$ is $(2\lambda L_g)$-gradient Lipschitz and $(\lambda \mu /2)$-strongly convex in $y$. This shows that if $(x,y)$ is an $\epsilon$-stationary point of $\fL_{\lambda}(x,y)$, then $x$ is also an $\fO(\kappa \epsilon)$-stationary point of $\varphi(x)$.
\textbf{Conversely}, if $x$ is an $\epsilon$-stationary point of $\varphi(x)$, the it suffices to find $y$ such that $ \Vert \nabla_y f(x,y) \Vert \le \fO(\epsilon/\kappa)$, then $(x,y)$ is a $2\epsilon$-stationary point of $\fL_{\lambda}(x,y)$. Under the deterministic setting, this goal can be achieved with $\fO(\kappa \log(\nicefrac{\kappa}{\epsilon}))$ first-order oracle calls by Theorem \ref{thm:GD-SC}. Under the stochastic setting, this goal can be achieved with $\fO(\kappa^2 (\sigma_f^2 + \lambda^2 \sigma_g^2)  \epsilon^{-2} )$ first-order oracle calls by \citep[Theorem 3]{allen2018make}.
\end{proof}

Note that one has to pay an additional $\fO(\kappa)$ factor when translating a stationary point of $\fL_{\lambda}(x,y)$ using the notation of \citet{shen2023penalty} to that of $\varphi(x)$ using our notation. Therefore, our notation is stronger than that of \citet{shen2023penalty}.

\subsection{Discussion on the Additional Assumptions in \citep{kwon2023fully}}  \label{apx:diss-kwon}
F${}^2$SA ~\citep[Algorithm 1]{kwon2023fully} additionally uses the following assumption.

\begin{asm} \label{asm:F2SA}[Assumption 4 and 5 in \citet{kwon2023fully}]
Besides Assumption \ref{asm:sc}, they also assume that
\begin{enumerate}
    \item $f(x,y)$ is $C_f$-Lipschitz in $x$;
    \item $g(x,y)$ is $C_g$-Lipschitz in $y$;
    \item $f(x,y)$ is $\rho_f$-Hessian Lipschitz in $(x,y)$.
\end{enumerate}
\end{asm}

The role of Assumption \ref{asm:F2SA} is to ensure the Lipschitz continuity of $y_{\lambda}^*(x)$ in $\lambda$ such that the algorithm can use an increasing $\lambda$. 
However, this additional assumption is unnecessary if we use a fixed $\lambda$. \citet{kwon2023fully} sets $\lambda_{t+1} = \lambda_t + \delta_t$ with the requirement 
\begin{align*}
    \frac{\delta_t}{\lambda_t} \le \frac{1}{16} \min \left\{1, \frac{T \mu}{L_g} \right\}.
\end{align*}
See \citep[Theorem 4, Condition (3b)]{kwon2023fully}. Therefore, setting $\delta_t = 0$ also meets the requirement.
In this case, when we discuss the algorithm by \citet{kwon2023fully}, we refer to this version with a fixed $\lambda$. This allows a fair comparison with our algorithm. After both algorithms use the same penalty $\lambda$, the only difference lies in our step size, which is the key to our improved analysis.

% To translate the notion of stationarity based on ∇f to our notion of stationarity, we need to pay an
% additional factor of O(κ log(1/ǫ)) or O(ǫ−2) in the two settings. In this sense, our notion of stationarity
% is stronger than the notion based on ∇f in the literature;



% \begin{remark}
% Similarly, we can also implement perturbed {F${}^2$BA} (Algorithm \ref{alg:PF2BA}) under the distributed setting. This would lead to a distributed first-order algorithm that can provably find an $\epsilon$-second-order stationary point with $\tilde \fO(\epsilon^{-2})$ complexity.
% \end{remark}


% \subsection{Numerical Experiments}

% \begin{figure*}[t]
%     \centering
%     \includegraphics[scale= 0.3]{fig/l2reg.eps} 
%     \caption{ Experiments of distributed learnable regularization on the 20news dataset.
%     }
%     \label{fig:l2reg}
% \end{figure*}

% We also use bilevel optimization to find the optimal regularization coefficient on the 20news dataset as~\citep{ji2021bilevel,liu2022bome}. In this formulation, the upper and lower-level functions take the form of:
% \begin{align*}
%     f(x,y) &:= \frac{1}{\vert \fD^{\rm val} \vert} \sum_{(a_i,b_i) \in \fD^{\rm val}} \ell(\, \langle a_i,y \rangle\,,b_i), \\g(x,y) &:= \frac{1}{\vert \fD^{\rm tr} \vert} \sum_{(a_i,b_i) \in \fD^{\rm tr}} 
% \ell \left( \,\langle a_i,y \rangle\,,b_i \right) + y^\top \Sigma(x) \,y,
% \end{align*}
% where $\ell(\,\cdot\,, \,\cdot\,)$ is the cross-entropy loss and $\Sigma(x) := {\rm diag}(\exp(x))$ determines the weights of features. We evenly distribute the datasets 
% on four agents and use the Gloo backend to implement parallel computing. We compare the distributed F${}^2$BA (Algorithm \ref{alg:CF2BA}) with AggITD~\citep{xiao2023communication} and AID with HIGD oracle~\citep{chen2022decentralized}. 
% The hyperparameter selection strategy follows the same way as Section \ref{sec:exp-clean}. 
% %For all the algorithms, we set the number of inner loops as 10 and tune the learning rates in $ \{1000,100,10,1,0.1,0.01,0.001 \}$.  For F${}^2$BA, we additionally tune the multiplier $\lambda$ in $\{10,100,1000\}$.
% We run all the algorithms for 5,000 epochs and plot the testing loss in Figure \ref{fig:l2reg}. It can be seen that the F${}^2$BA is more efficient than second-order methods in both computation and communication.

% \section{Proofs of Distributed Algorithms}

% First of all, we recall
% the convergence results of gradient tracking methods ~\cite{qu2017harnessing} for strongly convex objectives.

% \begin{prop}
% Suppose that Assumption \ref{asm:W} holds.
% Consider optimizing
% \begin{align*}
%     \min_{x \in \BR^d} \varphi(x):= \frac{1}{m} \sum_{i=1}^m f_i(x),
% \end{align*}
% where each $f_i(x)$ is  $L$-smooth and $\mu$-strongly convex.
% The update 
% \begin{align*} 
% X_{t+1} &= W^S \cdot (X_t - \eta G_t) \\
% G_{t+1} &= W^S \cdot (G_t + \nabla f(X_{t+1}) - \nabla f(X_t))
% \end{align*}
% with $\eta = 1/L, ~ S = \lceil \chi \log (10\kappa)\rceil$
% can find
% \begin{align*}
%     \sqrt{m }\Vert \bar x_{t} - x^* \Vert + \Vert X_{t} - \vone \bar x_{t} \Vert_F + \eta \Vert G_{t} - \vone \bar g_{t} \Vert_F \le \zeta
% \end{align*}
% within an iteration number bounded by
% \begin{align*}
%     \fO \left( 
%     \kappa \log \left( \dfrac{\sqrt{m}\Vert \bar x_0- x^* \Vert +\Vert X_0 - \vone \bar x_{0} \Vert_F +\eta \Vert G_{0} - \vone \bar g_{0} \Vert_F }{\zeta} \right)
%     \right).
% \end{align*}
% \end{prop}

% \begin{proof}
% We begin from
% \begin{align*}
% &\quad \Vert \bar x_{t+1} - x^* \Vert \\
% &=  \Vert \bar x_{t+1} - \eta \bar g_t - x^* \Vert \\
% &\le \Vert \bar x_{t+1} - \eta  \nabla f( \bar x_t) - x^* \Vert + \eta \Vert \nabla f(\bar x_t) -  \bar g_t\Vert\\
% &\le (1 - \mu/L)\Vert \bar x_t - x^* \Vert + \frac{\eta L}{\sqrt{m}} \Vert X_t - \vone \bar x_t \Vert_F, 
% \end{align*}
% where in the last step we use Lemma 10 and Lemma 8(c) in \cite{qu2017harnessing}.

% Also, it holds that
% \begin{align} \label{eq:cons}
% \Vert X_{t+1} - \vone \bar x_{t+1} \Vert \le \theta^S \Vert X_t - \vone \bar x_t \Vert + \theta^S \eta \Vert G_t - \vone \bar g_t \Vert. 
% \end{align}
% Furthermore,
% \begin{align*}
% &~\quad \Vert G_{t+1} - \vone \bar g_{t+1} \Vert_F \\
% &~\le \theta^S \Vert G_t - \vone \bar g_t \Vert_F + \theta^S \Vert \nabla f(X_{t+1}) - \nabla f(X_t) \Vert_F \\
% &~\le \theta^S \Vert G_t - \vone \bar g_t \Vert_F + \rho L \Vert X_{t+1} - X_t \Vert_F \\
% &~\le \theta^S \Vert G_t -  \vone \bar g_t \Vert_F + 3 \theta^S L \Vert X_{t} - \vone \bar x_t \Vert_F \\
% &~\quad + 3 \theta^S L \Vert \vone \bar x_{t+1} - \vone \bar x_t \Vert_F + 3\theta^S L \Vert X_{t+1} - \vone \bar x_{t+1} \Vert_F \\
% &\overset{(\ref{eq:cons})}{\le} \rho(1+ 3 \theta^S L \eta) \Vert G_t - \vone \bar g_t \Vert_F + 4 \theta^S L \Vert X_t - \vone \bar  x_t \Vert_F + 3 \theta^S L  \sqrt{m} \Vert
% \bar x_{t+1} - \bar x_t \Vert  \\
% &~\le~ \rho(1+ 3 \theta^S L \eta) \Vert G_t -  \vone \bar g_t \Vert_F + 4 \theta^S L \Vert X_t - \vone \bar x_t \Vert_F \\
% &~\quad +3 \theta^S L \eta \sqrt{m} \Vert  \nabla f(\bar x_t ) - \bar g_t \Vert + 3 \theta^S L \eta \sqrt{m} \Vert \nabla f(\bar x_t) \Vert \\
% &~\le  ~\theta^S(1+ 3 \theta^S L \eta) \Vert G_t -  \vone \bar g_t \Vert_F + \theta^S (4 L + 3  L^2 \eta) \Vert X_t - \vone \bar x_t \Vert_F   + 3 \theta^S L^2 \eta \sqrt{m} \Vert \bar x_t - x^* \Vert,
% \end{align*}
% where in the last step we use Lemma 8(c) in \cite{qu2017harnessing}. Combing the two inequalities above, we obtain
% \begin{align*}
% &\quad \Vert X_{t+1} - \vone \bar x_{t+1} \Vert_F + \eta \Vert G_{t+1} - \vone \bar g_{t+1} \Vert_F \\
% &\le \theta^S (1 + 4 L \eta + 3 L^2 \eta^2) \Vert X_t - \vone \bar x_t \Vert_F + \theta^S  (1 + 3 \rho L \eta) \eta \Vert G_t -  \vone \bar g_t \Vert_F + 3 \theta^S L^2\eta^2 \sqrt{m} \Vert \bar x_t - x^* \Vert. 
% \end{align*}
% Combing the above inequalities, we obtain
% \begin{align*}
% &\quad \sqrt{m} \Vert \bar x_{t+1} - x^* \Vert + \Vert X_{t+1} - \vone \bar x_{t+1} \Vert_F + \eta \Vert G_{t+1} - \vone \bar g_{t+1} \Vert_F \\
% &\le (1 - \mu/ L +  3 \theta^S L^2 \eta^2   ) \sqrt{m}\Vert \bar x_t - x^* \Vert + \theta^S (1+ 5L \eta + 3 L^2 \eta^2) \Vert X_t - \vone X_t \Vert_F +  \theta^S  (1 + 3 \theta^S L \eta) \eta \Vert G_t -  \vone \bar g_t \Vert_F.
% \end{align*}
% It can be verified the choice of $S$ ensures
% \begin{align*}
%     \theta^S \le \min \left \{ \frac{\mu}{6L}, ~ \frac{1}{10} \right\},
% \end{align*}
% then we can prove the following linear convergence:
% \begin{align*}
% &\quad   \sqrt{m} \Vert \bar x_{t+1} - x^* \Vert + \Vert X_{t+1} - \vone \bar x_{t+1} \Vert_F + \eta \Vert G_{t+1} - \vone \bar g_{t+1} \Vert_F \\
% &\le \left(1-\theta^S \right) \left( \sqrt{m}\Vert \bar x_t - x^* \Vert +  \Vert X_t - \vone X_t \Vert_F +  \eta \Vert G_t -  \vone \bar g_t \Vert_F \right).
% \end{align*}
% We arrive at the conclusion by telescoping the above inequality.
% \end{proof}

% \propGTGD*

% \begin{proof}
% In our proof, we define $\zeta: = \Vert \hat \nabla \fL_{\lambda}^*(x) - \nabla \fL_{\lambda}^*(x) \Vert_F$. The descent lemma tells us that
% \begin{align*}
% \fL_{\lambda}^*(\bar x_{t+1}) &\le \fL_{\lambda}^*(\bar x_t) + \langle \nabla \fL_{\lambda}^*(\bar x_t), \bar x_{t+1} - \bar x_t \rangle + \frac{L}{2}
% \Vert \bar x_{t+1} - \bar x_t \Vert^2 \\
% &= \fL_{\lambda}^*(\bar x_t) - \frac{\eta}{2} \Vert \nabla \fL_{\lambda}^*(\bar x_t) \Vert^2 -
% \left( \frac{\eta}{2} - \frac{L \eta^2}{2} \right) \Vert \bar g_t \Vert^2 + \frac{\eta}{2} \Vert \nabla \fL_{\lambda}^*(\bar x_t) - \bar g_t \Vert^2 \\
% &\le  \fL_{\lambda}^*(\bar x_t) - \frac{\eta}{2} \Vert \nabla \fL_{\lambda}^*(\bar x_t) \Vert^2 -
% \left( \frac{\eta}{2} - \frac{L \eta^2}{2} \right) \Vert \bar g_t \Vert^2 \\
% &\quad+ \eta \left\Vert \nabla \fL_{\lambda}^*(\bar x_t) - \frac{1}{m} \mathbf{1}^\top \nabla \fL_{\lambda}^*(X_t) \right\Vert^2
% + \frac{\eta}{m} \left \Vert  \hat \nabla \fL_{\lambda}^*(X_t) - \nabla \fL_{\lambda}^*(X_t)  \right \Vert_F^2. 
% \end{align*}
% Additionally,
% \begin{align*}
% &\quad \Vert G_{t+1} - \vone \bar g_{t+1} \Vert \\
% &\le \theta^{S_{\rm out}} \Vert G_t - \vone \bar g_t \Vert + \theta^{S_{\rm out}} \Vert \hat \nabla \fL_{\lambda}^*(X_{t+1}) - \hat \nabla \fL_{\lambda}^*(X_t) \Vert_F \\
% &\le \theta^{S_{\rm out}} \Vert G_t - \vone \bar g_t \Vert_F + \theta^{S_{\rm out}} L \Vert X_{t+1} - X_t \Vert_F + 2 \theta^{S_{\rm out}}  \zeta \\
% &\le \theta^{S_{\rm out}} \Vert G_t -  \vone \bar g_t \Vert_F + 3 \theta^{S_{\rm out}} L \Vert X_{t} - \vone \bar x_t \Vert_F  \\
% &\quad + 3 \theta^{S_{\rm out}} L \sqrt{m}\Vert  \bar x_{t+1} -  \bar x_t \Vert + 3\theta^{S_{\rm out}} L \Vert X_{t+1} - \vone \bar x_{t+1} \Vert_F + 2 \theta^{S_{\rm out}}  \zeta \\
% &\le \theta^{S_{\rm out}}(1+ 3 \theta^{S_{\rm out}} L \eta) \Vert G_t -  \vone \bar g_t \Vert_F + 4 \theta^{S_{\rm out}} L \Vert X_t - \vone \bar x_t \Vert_F + 3 \theta^{S_{\rm out}}   \eta \sqrt{m} \Vert \bar  g_{t} \Vert +2 \theta^{S_{\rm out}} \zeta.
% \end{align*}
% Combing this inequality with (\ref{eq:cons}) yields
% \begin{align*}
% &\quad \Vert X_{t+1} - \vone \bar x_{t+1} \Vert_F + \eta \Vert G_{t+1} - \vone \bar g_{t+1} \Vert_F \\
% &\le \rho (1 + 4 L \eta ) \Vert X_t - \vone \bar x_t \Vert_F + \theta^{S_{\rm out}}  (1 + 3 \theta^{S_{\rm out}} L \eta) \eta \Vert G_t -  \vone \bar g_t \Vert_F + 3 \theta^{S_{\rm out}} L \eta^2 \sqrt{m} \Vert \bar g_t \Vert + 2 \theta^{S_{\rm out}} \eta \zeta. 
% \end{align*}
% Taking square,
% \begin{align*}
%     &\quad \Vert X_{t+1} - \vone \bar x_{t+1} \Vert_F^2 + \eta^2 \Vert G_{t+1} - \vone \bar g_{t+1} \Vert_F^2 \\
% &\le 4 \theta^{2S_{\rm out}} (1 + 4 L \eta)^2 \Vert X_t - \vone \bar x_t \Vert_F^2 + 4 \theta^{2S_{\rm out}} (1+ 3 \theta^{S_{\rm out}} L \eta)^2 \eta^2 \Vert G_t - \vone  \bar g_t \Vert_F^2 \\
% &\quad + 36 \theta^{2S_{\rm out}} L^2 \eta^4 m \Vert \bar g_t \Vert^2 + 16 \theta^{2S_{\rm out}} \eta^2  \zeta^2.
% \end{align*}
% Therefore,
% \begin{align*}
%     &\quad \fL_{\lambda}^*(\bar x_{t+1})  + \frac{1}{m \eta } \Vert X_{t+1} - \vone \bar x_{t+1} \Vert_F^2 + \frac{\eta}{m } \Vert G_{t+1} - \vone \bar g_{t+1} \Vert_F^2 \\
%     &\le \fL_{\lambda}^*(\bar x_t) - \frac{\eta}{2} \Vert \nabla \fL_{\lambda}^*(\bar x_t) \Vert^2 - \left( \frac{\eta}{2} - \frac{L \eta^2}{2} - 36 \theta^{2S_{\rm out}} L^2 \eta^3  \right) \Vert \bar g_t \Vert^2 \\
%     &\quad +
% \end{align*}
% \end{proof}
